%% verse_protrusion.tex
%% Microtype protrusion for verse numbers (edition-portable).
%% Requires, loaded BEFORE this file: fontspec (main font set), microtype, scripture.
%% Effect: a verse number at the start of a line hangs a tuned partial amount
%% into the left margin. Mid-line numbers are unaffected (microtype only
%% protrudes line-edge glyphs). Font-agnostic: uses OpenType superior figures
%% of the CURRENT body font; falls back to \textsuperscript if unavailable.
%%
%% WORKING MECHANISM (the bits that are easy to get wrong):
%%   1. *** NO named context, NO \microtypecontext. *** Use a PLAIN, default
%%      \SetProtrusion (neither [name=...] nor [context=...]). microtype then
%%      MERGES these entries into whatever protrusion list it loads for any
%%      matching font, so the superior-figure instance gets the codepoints
%%      automatically. The earlier [context=verseprot] +
%%      \microtypecontext{protrusion=verseprot} approach was a NO-OP: the
%%      verbose log showed the context switch loaded an auto-generated EMPTY
%%      list named `<file>/<line>' and the named list never bound to the active
%%      superior font, so lpcode stayed 0 (measured: zero movement at 600 dpi).
%%      The global config is safe here because the superscript-codepoint glyphs
%%      (U+00B9 etc.) appear in this document ONLY as verse numbers (body prose
%%      never emits superior figures), so the entries are effectively scoped to
%%      verse numbers already. PROOF the binding now works: with the global
%%      config, `\the\lpcode\font 185' (U+00B9) reports a nonzero value on the
%%      VerticalPosition=Superior instance (was 0/unset before), and 600 dpi
%%      measurement shows the leading verse number's leftmost ink ~9 px left of
%%      where it sits with the feature off.
%%   2. Font set is the full 5-axis wildcard { font = */*/*/*/* } so it matches
%%      BOTH the body font instance AND the separate font instance that fontspec
%%      spins up for VerticalPosition=Superior (e.g. EBGaramond(1)).
%%   3. Each entry is char = {left,right} (ONE brace, comma-separated). Left-only
%%      hang -> {350,0}. Two separate braces {350}{0} is wrong syntax.
%%   4. To retune the factor, edit \verseprot@factor (one place). The \edef
%%      below expands it into all ten entries before passing to \SetProtrusion.
%%      TWO critical constraints on the expansion approach:
%%      (a) microtype stores \SetProtrusion's third arg via \gdef -- if we pass
%%          `{\verseprot@list}' as the arg, the macro TOKEN (not its expansion)
%%          is stored, and microtype's \MT@set@codes later sees \verseprot@list
%%          as a single opaque token, not as the hex-entry list. We therefore use
%%          `\edef\verseprot@do@protrusion{...{\verseprot@list}...}' to EXPAND the
%%          list inline and store the literal tokens in \verseprot@do@protrusion
%%          before \SetProtrusion is ever called.
%%      (b) microtype reads the \SetProtrusion list with \catcode` '=9 (spaces
%%          ignored), but our \edef runs with normal catcodes where space is
%%          catcode 10. If the expanded list contains catcode-10 space tokens,
%%          \MT@pr@split@val (which reads `left,right\relax') sees a leading space
%%          before `{350,0}' and can't find the `,\relax' terminator -> runaway.
%%          The fix: write the list with NO spaces (all newlines suppressed by `%'
%%          and no space around `=' or `{'). microtype's own parser ignores the
%%          lack of spaces; it just needs the correct token sequence.
%%   5. *** THE OTHER KEY FIX ***  microtype (luatex) keys protrusion by the codepoint
%%      of the glyph AS IT APPEARS IN THE NODE LIST, i.e. AFTER GSUB shaping.
%%      VerticalPosition=Superior runs the OpenType `sups' feature, which in
%%      EB Garamond substitutes each ASCII digit for the glyph at its Unicode
%%      SUPERSCRIPT codepoint:
%%          0->U+2070  1->U+00B9  2->U+00B2  3->U+00B3
%%          4..9 -> U+2074..U+2079
%%      Therefore the list MUST be keyed on those superscript codepoints, NOT on
%%      the ASCII digits 0-9. (Keying 0-9 = U+0030..U+0039 was the bug: those
%%      input glyphs never reach the line because sups replaces them, so the
%%      lpcode never bound to the rendered superior glyph -> ZERO protrusion,
%%      verified by 600 dpi pixel measurement and by dumping the substituted
%%      glyph names with fontTools: three -> uni00B3, etc.)
%%   6. microtype protrudes GLYPH nodes, not boxes. The substituted superior
%%      figures are real glyphs and protrude; \textsuperscript (a raised box)
%%      does not -- hence the OpenType path is required for the hang, with
%%      \textsuperscript only as the non-protruding fallback when the font lacks
%%      superior figures.
%% Verified by measurement (tests/verse_protrusion_visual_check.tex, 600 dpi):
%% with this setup the leading verse number's leftmost ink sits measurably LEFT
%% of the column rule, while mid-line numbers stay inline; feature ON vs OFF
%% now differs at the left edge.
%%
%% PORTABILITY NOTE: the superscript-codepoint mapping above is what EB Garamond
%% (and most fonts following the Unicode/AGL convention) emit for `sups'. A font
%% whose superior figures live at private-use or `.sups'-named glyph slots with
%% no Unicode value would need different keys; in that case microtype cannot
%% address them by codepoint and the \textsuperscript fallback (no protrusion)
%% is the safe outcome.
\makeatletter

% --- Precondition guards: raise a clear error if required packages are absent.
%     Without microtype (\SetProtrusion) or fontspec (\addfontfeature) the file
%     would fail later with a cryptic "Undefined control sequence" inside an \edef.
\@ifundefined{SetProtrusion}{%
  \PackageError{verse_protrusion}%
    {microtype must be loaded before verse_protrusion.tex}%
    {Add \string\usepackage{microtype} (and fontspec) before \string%% verse_protrusion.tex
%% Microtype protrusion for verse numbers (edition-portable).
%% Requires, loaded BEFORE this file: fontspec (main font set), microtype, scripture.
%% Effect: a verse number at the start of a line hangs a tuned partial amount
%% into the left margin. Mid-line numbers are unaffected (microtype only
%% protrudes line-edge glyphs). Font-agnostic: uses OpenType superior figures
%% of the CURRENT body font; falls back to \textsuperscript if unavailable.
%%
%% WORKING MECHANISM (the bits that are easy to get wrong):
%%   1. *** NO named context, NO \microtypecontext. *** Use a PLAIN, default
%%      \SetProtrusion (neither [name=...] nor [context=...]). microtype then
%%      MERGES these entries into whatever protrusion list it loads for any
%%      matching font, so the superior-figure instance gets the codepoints
%%      automatically. The earlier [context=verseprot] +
%%      \microtypecontext{protrusion=verseprot} approach was a NO-OP: the
%%      verbose log showed the context switch loaded an auto-generated EMPTY
%%      list named `<file>/<line>' and the named list never bound to the active
%%      superior font, so lpcode stayed 0 (measured: zero movement at 600 dpi).
%%      The global config is safe here because the superscript-codepoint glyphs
%%      (U+00B9 etc.) appear in this document ONLY as verse numbers (body prose
%%      never emits superior figures), so the entries are effectively scoped to
%%      verse numbers already. PROOF the binding now works: with the global
%%      config, `\the\lpcode\font 185' (U+00B9) reports a nonzero value on the
%%      VerticalPosition=Superior instance (was 0/unset before), and 600 dpi
%%      measurement shows the leading verse number's leftmost ink ~9 px left of
%%      where it sits with the feature off.
%%   2. Font set is the full 5-axis wildcard { font = */*/*/*/* } so it matches
%%      BOTH the body font instance AND the separate font instance that fontspec
%%      spins up for VerticalPosition=Superior (e.g. EBGaramond(1)).
%%   3. Each entry is char = {left,right} (ONE brace, comma-separated). Left-only
%%      hang -> {350,0}. Two separate braces {350}{0} is wrong syntax.
%%   4. To retune the factor, edit \verseprot@factor (one place). The \edef
%%      below expands it into all ten entries before passing to \SetProtrusion.
%%      TWO critical constraints on the expansion approach:
%%      (a) microtype stores \SetProtrusion's third arg via \gdef -- if we pass
%%          `{\verseprot@list}' as the arg, the macro TOKEN (not its expansion)
%%          is stored, and microtype's \MT@set@codes later sees \verseprot@list
%%          as a single opaque token, not as the hex-entry list. We therefore use
%%          `\edef\verseprot@do@protrusion{...{\verseprot@list}...}' to EXPAND the
%%          list inline and store the literal tokens in \verseprot@do@protrusion
%%          before \SetProtrusion is ever called.
%%      (b) microtype reads the \SetProtrusion list with \catcode` '=9 (spaces
%%          ignored), but our \edef runs with normal catcodes where space is
%%          catcode 10. If the expanded list contains catcode-10 space tokens,
%%          \MT@pr@split@val (which reads `left,right\relax') sees a leading space
%%          before `{350,0}' and can't find the `,\relax' terminator -> runaway.
%%          The fix: write the list with NO spaces (all newlines suppressed by `%'
%%          and no space around `=' or `{'). microtype's own parser ignores the
%%          lack of spaces; it just needs the correct token sequence.
%%   5. *** THE OTHER KEY FIX ***  microtype (luatex) keys protrusion by the codepoint
%%      of the glyph AS IT APPEARS IN THE NODE LIST, i.e. AFTER GSUB shaping.
%%      VerticalPosition=Superior runs the OpenType `sups' feature, which in
%%      EB Garamond substitutes each ASCII digit for the glyph at its Unicode
%%      SUPERSCRIPT codepoint:
%%          0->U+2070  1->U+00B9  2->U+00B2  3->U+00B3
%%          4..9 -> U+2074..U+2079
%%      Therefore the list MUST be keyed on those superscript codepoints, NOT on
%%      the ASCII digits 0-9. (Keying 0-9 = U+0030..U+0039 was the bug: those
%%      input glyphs never reach the line because sups replaces them, so the
%%      lpcode never bound to the rendered superior glyph -> ZERO protrusion,
%%      verified by 600 dpi pixel measurement and by dumping the substituted
%%      glyph names with fontTools: three -> uni00B3, etc.)
%%   6. microtype protrudes GLYPH nodes, not boxes. The substituted superior
%%      figures are real glyphs and protrude; \textsuperscript (a raised box)
%%      does not -- hence the OpenType path is required for the hang, with
%%      \textsuperscript only as the non-protruding fallback when the font lacks
%%      superior figures.
%% Verified by measurement (tests/verse_protrusion_visual_check.tex, 600 dpi):
%% with this setup the leading verse number's leftmost ink sits measurably LEFT
%% of the column rule, while mid-line numbers stay inline; feature ON vs OFF
%% now differs at the left edge.
%%
%% PORTABILITY NOTE: the superscript-codepoint mapping above is what EB Garamond
%% (and most fonts following the Unicode/AGL convention) emit for `sups'. A font
%% whose superior figures live at private-use or `.sups'-named glyph slots with
%% no Unicode value would need different keys; in that case microtype cannot
%% address them by codepoint and the \textsuperscript fallback (no protrusion)
%% is the safe outcome.
\makeatletter

% --- Precondition guards: raise a clear error if required packages are absent.
%     Without microtype (\SetProtrusion) or fontspec (\addfontfeature) the file
%     would fail later with a cryptic "Undefined control sequence" inside an \edef.
\@ifundefined{SetProtrusion}{%
  \PackageError{verse_protrusion}%
    {microtype must be loaded before verse_protrusion.tex}%
    {Add \string\usepackage{microtype} (and fontspec) before \string%% verse_protrusion.tex
%% Microtype protrusion for verse numbers (edition-portable).
%% Requires, loaded BEFORE this file: fontspec (main font set), microtype, scripture.
%% Effect: a verse number at the start of a line hangs a tuned partial amount
%% into the left margin. Mid-line numbers are unaffected (microtype only
%% protrudes line-edge glyphs). Font-agnostic: uses OpenType superior figures
%% of the CURRENT body font; falls back to \textsuperscript if unavailable.
%%
%% WORKING MECHANISM (the bits that are easy to get wrong):
%%   1. *** NO named context, NO \microtypecontext. *** Use a PLAIN, default
%%      \SetProtrusion (neither [name=...] nor [context=...]). microtype then
%%      MERGES these entries into whatever protrusion list it loads for any
%%      matching font, so the superior-figure instance gets the codepoints
%%      automatically. The earlier [context=verseprot] +
%%      \microtypecontext{protrusion=verseprot} approach was a NO-OP: the
%%      verbose log showed the context switch loaded an auto-generated EMPTY
%%      list named `<file>/<line>' and the named list never bound to the active
%%      superior font, so lpcode stayed 0 (measured: zero movement at 600 dpi).
%%      The global config is safe here because the superscript-codepoint glyphs
%%      (U+00B9 etc.) appear in this document ONLY as verse numbers (body prose
%%      never emits superior figures), so the entries are effectively scoped to
%%      verse numbers already. PROOF the binding now works: with the global
%%      config, `\the\lpcode\font 185' (U+00B9) reports a nonzero value on the
%%      VerticalPosition=Superior instance (was 0/unset before), and 600 dpi
%%      measurement shows the leading verse number's leftmost ink ~9 px left of
%%      where it sits with the feature off.
%%   2. Font set is the full 5-axis wildcard { font = */*/*/*/* } so it matches
%%      BOTH the body font instance AND the separate font instance that fontspec
%%      spins up for VerticalPosition=Superior (e.g. EBGaramond(1)).
%%   3. Each entry is char = {left,right} (ONE brace, comma-separated). Left-only
%%      hang -> {350,0}. Two separate braces {350}{0} is wrong syntax.
%%   4. To retune the factor, edit \verseprot@factor (one place). The \edef
%%      below expands it into all ten entries before passing to \SetProtrusion.
%%      TWO critical constraints on the expansion approach:
%%      (a) microtype stores \SetProtrusion's third arg via \gdef -- if we pass
%%          `{\verseprot@list}' as the arg, the macro TOKEN (not its expansion)
%%          is stored, and microtype's \MT@set@codes later sees \verseprot@list
%%          as a single opaque token, not as the hex-entry list. We therefore use
%%          `\edef\verseprot@do@protrusion{...{\verseprot@list}...}' to EXPAND the
%%          list inline and store the literal tokens in \verseprot@do@protrusion
%%          before \SetProtrusion is ever called.
%%      (b) microtype reads the \SetProtrusion list with \catcode` '=9 (spaces
%%          ignored), but our \edef runs with normal catcodes where space is
%%          catcode 10. If the expanded list contains catcode-10 space tokens,
%%          \MT@pr@split@val (which reads `left,right\relax') sees a leading space
%%          before `{350,0}' and can't find the `,\relax' terminator -> runaway.
%%          The fix: write the list with NO spaces (all newlines suppressed by `%'
%%          and no space around `=' or `{'). microtype's own parser ignores the
%%          lack of spaces; it just needs the correct token sequence.
%%   5. *** THE OTHER KEY FIX ***  microtype (luatex) keys protrusion by the codepoint
%%      of the glyph AS IT APPEARS IN THE NODE LIST, i.e. AFTER GSUB shaping.
%%      VerticalPosition=Superior runs the OpenType `sups' feature, which in
%%      EB Garamond substitutes each ASCII digit for the glyph at its Unicode
%%      SUPERSCRIPT codepoint:
%%          0->U+2070  1->U+00B9  2->U+00B2  3->U+00B3
%%          4..9 -> U+2074..U+2079
%%      Therefore the list MUST be keyed on those superscript codepoints, NOT on
%%      the ASCII digits 0-9. (Keying 0-9 = U+0030..U+0039 was the bug: those
%%      input glyphs never reach the line because sups replaces them, so the
%%      lpcode never bound to the rendered superior glyph -> ZERO protrusion,
%%      verified by 600 dpi pixel measurement and by dumping the substituted
%%      glyph names with fontTools: three -> uni00B3, etc.)
%%   6. microtype protrudes GLYPH nodes, not boxes. The substituted superior
%%      figures are real glyphs and protrude; \textsuperscript (a raised box)
%%      does not -- hence the OpenType path is required for the hang, with
%%      \textsuperscript only as the non-protruding fallback when the font lacks
%%      superior figures.
%% Verified by measurement (tests/verse_protrusion_visual_check.tex, 600 dpi):
%% with this setup the leading verse number's leftmost ink sits measurably LEFT
%% of the column rule, while mid-line numbers stay inline; feature ON vs OFF
%% now differs at the left edge.
%%
%% PORTABILITY NOTE: the superscript-codepoint mapping above is what EB Garamond
%% (and most fonts following the Unicode/AGL convention) emit for `sups'. A font
%% whose superior figures live at private-use or `.sups'-named glyph slots with
%% no Unicode value would need different keys; in that case microtype cannot
%% address them by codepoint and the \textsuperscript fallback (no protrusion)
%% is the safe outcome.
\makeatletter

% --- Precondition guards: raise a clear error if required packages are absent.
%     Without microtype (\SetProtrusion) or fontspec (\addfontfeature) the file
%     would fail later with a cryptic "Undefined control sequence" inside an \edef.
\@ifundefined{SetProtrusion}{%
  \PackageError{verse_protrusion}%
    {microtype must be loaded before verse_protrusion.tex}%
    {Add \string\usepackage{microtype} (and fontspec) before \string%% verse_protrusion.tex
%% Microtype protrusion for verse numbers (edition-portable).
%% Requires, loaded BEFORE this file: fontspec (main font set), microtype, scripture.
%% Effect: a verse number at the start of a line hangs a tuned partial amount
%% into the left margin. Mid-line numbers are unaffected (microtype only
%% protrudes line-edge glyphs). Font-agnostic: uses OpenType superior figures
%% of the CURRENT body font; falls back to \textsuperscript if unavailable.
%%
%% WORKING MECHANISM (the bits that are easy to get wrong):
%%   1. *** NO named context, NO \microtypecontext. *** Use a PLAIN, default
%%      \SetProtrusion (neither [name=...] nor [context=...]). microtype then
%%      MERGES these entries into whatever protrusion list it loads for any
%%      matching font, so the superior-figure instance gets the codepoints
%%      automatically. The earlier [context=verseprot] +
%%      \microtypecontext{protrusion=verseprot} approach was a NO-OP: the
%%      verbose log showed the context switch loaded an auto-generated EMPTY
%%      list named `<file>/<line>' and the named list never bound to the active
%%      superior font, so lpcode stayed 0 (measured: zero movement at 600 dpi).
%%      The global config is safe here because the superscript-codepoint glyphs
%%      (U+00B9 etc.) appear in this document ONLY as verse numbers (body prose
%%      never emits superior figures), so the entries are effectively scoped to
%%      verse numbers already. PROOF the binding now works: with the global
%%      config, `\the\lpcode\font 185' (U+00B9) reports a nonzero value on the
%%      VerticalPosition=Superior instance (was 0/unset before), and 600 dpi
%%      measurement shows the leading verse number's leftmost ink ~9 px left of
%%      where it sits with the feature off.
%%   2. Font set is the full 5-axis wildcard { font = */*/*/*/* } so it matches
%%      BOTH the body font instance AND the separate font instance that fontspec
%%      spins up for VerticalPosition=Superior (e.g. EBGaramond(1)).
%%   3. Each entry is char = {left,right} (ONE brace, comma-separated). Left-only
%%      hang -> {350,0}. Two separate braces {350}{0} is wrong syntax.
%%   4. To retune the factor, edit \verseprot@factor (one place). The \edef
%%      below expands it into all ten entries before passing to \SetProtrusion.
%%      TWO critical constraints on the expansion approach:
%%      (a) microtype stores \SetProtrusion's third arg via \gdef -- if we pass
%%          `{\verseprot@list}' as the arg, the macro TOKEN (not its expansion)
%%          is stored, and microtype's \MT@set@codes later sees \verseprot@list
%%          as a single opaque token, not as the hex-entry list. We therefore use
%%          `\edef\verseprot@do@protrusion{...{\verseprot@list}...}' to EXPAND the
%%          list inline and store the literal tokens in \verseprot@do@protrusion
%%          before \SetProtrusion is ever called.
%%      (b) microtype reads the \SetProtrusion list with \catcode` '=9 (spaces
%%          ignored), but our \edef runs with normal catcodes where space is
%%          catcode 10. If the expanded list contains catcode-10 space tokens,
%%          \MT@pr@split@val (which reads `left,right\relax') sees a leading space
%%          before `{350,0}' and can't find the `,\relax' terminator -> runaway.
%%          The fix: write the list with NO spaces (all newlines suppressed by `%'
%%          and no space around `=' or `{'). microtype's own parser ignores the
%%          lack of spaces; it just needs the correct token sequence.
%%   5. *** THE OTHER KEY FIX ***  microtype (luatex) keys protrusion by the codepoint
%%      of the glyph AS IT APPEARS IN THE NODE LIST, i.e. AFTER GSUB shaping.
%%      VerticalPosition=Superior runs the OpenType `sups' feature, which in
%%      EB Garamond substitutes each ASCII digit for the glyph at its Unicode
%%      SUPERSCRIPT codepoint:
%%          0->U+2070  1->U+00B9  2->U+00B2  3->U+00B3
%%          4..9 -> U+2074..U+2079
%%      Therefore the list MUST be keyed on those superscript codepoints, NOT on
%%      the ASCII digits 0-9. (Keying 0-9 = U+0030..U+0039 was the bug: those
%%      input glyphs never reach the line because sups replaces them, so the
%%      lpcode never bound to the rendered superior glyph -> ZERO protrusion,
%%      verified by 600 dpi pixel measurement and by dumping the substituted
%%      glyph names with fontTools: three -> uni00B3, etc.)
%%   6. microtype protrudes GLYPH nodes, not boxes. The substituted superior
%%      figures are real glyphs and protrude; \textsuperscript (a raised box)
%%      does not -- hence the OpenType path is required for the hang, with
%%      \textsuperscript only as the non-protruding fallback when the font lacks
%%      superior figures.
%% Verified by measurement (tests/verse_protrusion_visual_check.tex, 600 dpi):
%% with this setup the leading verse number's leftmost ink sits measurably LEFT
%% of the column rule, while mid-line numbers stay inline; feature ON vs OFF
%% now differs at the left edge.
%%
%% PORTABILITY NOTE: the superscript-codepoint mapping above is what EB Garamond
%% (and most fonts following the Unicode/AGL convention) emit for `sups'. A font
%% whose superior figures live at private-use or `.sups'-named glyph slots with
%% no Unicode value would need different keys; in that case microtype cannot
%% address them by codepoint and the \textsuperscript fallback (no protrusion)
%% is the safe outcome.
\makeatletter

% --- Precondition guards: raise a clear error if required packages are absent.
%     Without microtype (\SetProtrusion) or fontspec (\addfontfeature) the file
%     would fail later with a cryptic "Undefined control sequence" inside an \edef.
\@ifundefined{SetProtrusion}{%
  \PackageError{verse_protrusion}%
    {microtype must be loaded before verse_protrusion.tex}%
    {Add \string\usepackage{microtype} (and fontspec) before \string\input{verse_protrusion}.}%
}{}
\@ifundefined{addfontfeature}{%
  \PackageError{verse_protrusion}%
    {fontspec must be loaded before verse_protrusion.tex}%
    {Add \string\usepackage{fontspec} (and set a main font) before \string\input{verse_protrusion}.}%
}{}

% --- Tunable: partial left-protrusion factor (out of 1000 = one glyph width).
%     Edit ONLY this line to retune the hang; the \edef below propagates it
%     into all ten digit entries automatically.
\def\verseprot@factor{350}

% Build the entry list with the factor expanded in, NO spaces anywhere
% (catcode-10 spaces in the stored token list break \MT@pr@split@val --
%  see header note #4b).  Every newline is suppressed by a trailing `%'.
\edef\verseprot@list{%
"2070={\verseprot@factor,0},%
"00B9={\verseprot@factor,0},%
"00B2={\verseprot@factor,0},%
"00B3={\verseprot@factor,0},%
"2074={\verseprot@factor,0},%
"2075={\verseprot@factor,0},%
"2076={\verseprot@factor,0},%
"2077={\verseprot@factor,0},%
"2078={\verseprot@factor,0},%
"2079={\verseprot@factor,0}}

% --- Manual override (an edition may force the feature off):
\newif\ifverseprotrusion  \verseprotrusiontrue

% --- Global (default) microtype protrusion: left-only hang on the SUPERSCRIPT
%     glyphs that the `sups' feature substitutes for the digits 0-9 (see header
%     note #5). Keyed on the substituted Unicode codepoints, NOT on ASCII 0-9.
%     NO [name=...] and NO [context=...]: this MERGES into the protrusion list
%     microtype loads for the VerticalPosition=Superior font instance, which is
%     what actually binds the codepoints (header note #1). Safe globally because
%     these superscript glyphs only ever appear as verse numbers in this doc.
%     \verseprot@list is passed via \edef so the literal text (not the macro
%     name) is baked into the call before microtype reads it (header note #4).
% [load=default]: INHERIT microtype's built-in protrusion (quotes, hyphens,
% dashes, periods -- the hanging punctuation) and ADD the superscript-figure
% codepoints on top.  Without load=, a global \SetProtrusion REPLACES the
% font's whole protrusion vector, silently disabling all default hanging
% punctuation across the edition.
\edef\verseprot@do@protrusion{%
  \noexpand\SetProtrusion
    [ load = default ]%
    { font = */*/*/*/* }%
    {\verseprot@list}%
}%
\verseprot@do@protrusion

% Adapted for LaTeXGen: called in scripture's font hook AFTER selecting
% scripturefont. The sermon main font is intentionally a different family.
\newif\ifverseprot@use
\newcommand{\detectverseprotrusion}{%
  \setbox0=\hbox{{\addfontfeature{VerticalPosition=Superior}0}}%
  \setbox2=\hbox{0}%
  \ifdim\wd0=\wd2
    \verseprot@usefalse
  \else
    \ifverseprotrusion \verseprot@usetrue \else \verseprot@usefalse \fi
  \fi
}

% --- The formatter wired into scripture's verse/format.
\NewDocumentCommand\verseprotrudenum{m}{%
  \ifverseprot@use
    {\addfontfeature{VerticalPosition=Superior}#1}%
  \else
    \textsuperscript{#1}%
  \fi
}

\makeatother
.}%
}{}
\@ifundefined{addfontfeature}{%
  \PackageError{verse_protrusion}%
    {fontspec must be loaded before verse_protrusion.tex}%
    {Add \string\usepackage{fontspec} (and set a main font) before \string%% verse_protrusion.tex
%% Microtype protrusion for verse numbers (edition-portable).
%% Requires, loaded BEFORE this file: fontspec (main font set), microtype, scripture.
%% Effect: a verse number at the start of a line hangs a tuned partial amount
%% into the left margin. Mid-line numbers are unaffected (microtype only
%% protrudes line-edge glyphs). Font-agnostic: uses OpenType superior figures
%% of the CURRENT body font; falls back to \textsuperscript if unavailable.
%%
%% WORKING MECHANISM (the bits that are easy to get wrong):
%%   1. *** NO named context, NO \microtypecontext. *** Use a PLAIN, default
%%      \SetProtrusion (neither [name=...] nor [context=...]). microtype then
%%      MERGES these entries into whatever protrusion list it loads for any
%%      matching font, so the superior-figure instance gets the codepoints
%%      automatically. The earlier [context=verseprot] +
%%      \microtypecontext{protrusion=verseprot} approach was a NO-OP: the
%%      verbose log showed the context switch loaded an auto-generated EMPTY
%%      list named `<file>/<line>' and the named list never bound to the active
%%      superior font, so lpcode stayed 0 (measured: zero movement at 600 dpi).
%%      The global config is safe here because the superscript-codepoint glyphs
%%      (U+00B9 etc.) appear in this document ONLY as verse numbers (body prose
%%      never emits superior figures), so the entries are effectively scoped to
%%      verse numbers already. PROOF the binding now works: with the global
%%      config, `\the\lpcode\font 185' (U+00B9) reports a nonzero value on the
%%      VerticalPosition=Superior instance (was 0/unset before), and 600 dpi
%%      measurement shows the leading verse number's leftmost ink ~9 px left of
%%      where it sits with the feature off.
%%   2. Font set is the full 5-axis wildcard { font = */*/*/*/* } so it matches
%%      BOTH the body font instance AND the separate font instance that fontspec
%%      spins up for VerticalPosition=Superior (e.g. EBGaramond(1)).
%%   3. Each entry is char = {left,right} (ONE brace, comma-separated). Left-only
%%      hang -> {350,0}. Two separate braces {350}{0} is wrong syntax.
%%   4. To retune the factor, edit \verseprot@factor (one place). The \edef
%%      below expands it into all ten entries before passing to \SetProtrusion.
%%      TWO critical constraints on the expansion approach:
%%      (a) microtype stores \SetProtrusion's third arg via \gdef -- if we pass
%%          `{\verseprot@list}' as the arg, the macro TOKEN (not its expansion)
%%          is stored, and microtype's \MT@set@codes later sees \verseprot@list
%%          as a single opaque token, not as the hex-entry list. We therefore use
%%          `\edef\verseprot@do@protrusion{...{\verseprot@list}...}' to EXPAND the
%%          list inline and store the literal tokens in \verseprot@do@protrusion
%%          before \SetProtrusion is ever called.
%%      (b) microtype reads the \SetProtrusion list with \catcode` '=9 (spaces
%%          ignored), but our \edef runs with normal catcodes where space is
%%          catcode 10. If the expanded list contains catcode-10 space tokens,
%%          \MT@pr@split@val (which reads `left,right\relax') sees a leading space
%%          before `{350,0}' and can't find the `,\relax' terminator -> runaway.
%%          The fix: write the list with NO spaces (all newlines suppressed by `%'
%%          and no space around `=' or `{'). microtype's own parser ignores the
%%          lack of spaces; it just needs the correct token sequence.
%%   5. *** THE OTHER KEY FIX ***  microtype (luatex) keys protrusion by the codepoint
%%      of the glyph AS IT APPEARS IN THE NODE LIST, i.e. AFTER GSUB shaping.
%%      VerticalPosition=Superior runs the OpenType `sups' feature, which in
%%      EB Garamond substitutes each ASCII digit for the glyph at its Unicode
%%      SUPERSCRIPT codepoint:
%%          0->U+2070  1->U+00B9  2->U+00B2  3->U+00B3
%%          4..9 -> U+2074..U+2079
%%      Therefore the list MUST be keyed on those superscript codepoints, NOT on
%%      the ASCII digits 0-9. (Keying 0-9 = U+0030..U+0039 was the bug: those
%%      input glyphs never reach the line because sups replaces them, so the
%%      lpcode never bound to the rendered superior glyph -> ZERO protrusion,
%%      verified by 600 dpi pixel measurement and by dumping the substituted
%%      glyph names with fontTools: three -> uni00B3, etc.)
%%   6. microtype protrudes GLYPH nodes, not boxes. The substituted superior
%%      figures are real glyphs and protrude; \textsuperscript (a raised box)
%%      does not -- hence the OpenType path is required for the hang, with
%%      \textsuperscript only as the non-protruding fallback when the font lacks
%%      superior figures.
%% Verified by measurement (tests/verse_protrusion_visual_check.tex, 600 dpi):
%% with this setup the leading verse number's leftmost ink sits measurably LEFT
%% of the column rule, while mid-line numbers stay inline; feature ON vs OFF
%% now differs at the left edge.
%%
%% PORTABILITY NOTE: the superscript-codepoint mapping above is what EB Garamond
%% (and most fonts following the Unicode/AGL convention) emit for `sups'. A font
%% whose superior figures live at private-use or `.sups'-named glyph slots with
%% no Unicode value would need different keys; in that case microtype cannot
%% address them by codepoint and the \textsuperscript fallback (no protrusion)
%% is the safe outcome.
\makeatletter

% --- Precondition guards: raise a clear error if required packages are absent.
%     Without microtype (\SetProtrusion) or fontspec (\addfontfeature) the file
%     would fail later with a cryptic "Undefined control sequence" inside an \edef.
\@ifundefined{SetProtrusion}{%
  \PackageError{verse_protrusion}%
    {microtype must be loaded before verse_protrusion.tex}%
    {Add \string\usepackage{microtype} (and fontspec) before \string\input{verse_protrusion}.}%
}{}
\@ifundefined{addfontfeature}{%
  \PackageError{verse_protrusion}%
    {fontspec must be loaded before verse_protrusion.tex}%
    {Add \string\usepackage{fontspec} (and set a main font) before \string\input{verse_protrusion}.}%
}{}

% --- Tunable: partial left-protrusion factor (out of 1000 = one glyph width).
%     Edit ONLY this line to retune the hang; the \edef below propagates it
%     into all ten digit entries automatically.
\def\verseprot@factor{350}

% Build the entry list with the factor expanded in, NO spaces anywhere
% (catcode-10 spaces in the stored token list break \MT@pr@split@val --
%  see header note #4b).  Every newline is suppressed by a trailing `%'.
\edef\verseprot@list{%
"2070={\verseprot@factor,0},%
"00B9={\verseprot@factor,0},%
"00B2={\verseprot@factor,0},%
"00B3={\verseprot@factor,0},%
"2074={\verseprot@factor,0},%
"2075={\verseprot@factor,0},%
"2076={\verseprot@factor,0},%
"2077={\verseprot@factor,0},%
"2078={\verseprot@factor,0},%
"2079={\verseprot@factor,0}}

% --- Manual override (an edition may force the feature off):
\newif\ifverseprotrusion  \verseprotrusiontrue

% --- Global (default) microtype protrusion: left-only hang on the SUPERSCRIPT
%     glyphs that the `sups' feature substitutes for the digits 0-9 (see header
%     note #5). Keyed on the substituted Unicode codepoints, NOT on ASCII 0-9.
%     NO [name=...] and NO [context=...]: this MERGES into the protrusion list
%     microtype loads for the VerticalPosition=Superior font instance, which is
%     what actually binds the codepoints (header note #1). Safe globally because
%     these superscript glyphs only ever appear as verse numbers in this doc.
%     \verseprot@list is passed via \edef so the literal text (not the macro
%     name) is baked into the call before microtype reads it (header note #4).
% [load=default]: INHERIT microtype's built-in protrusion (quotes, hyphens,
% dashes, periods -- the hanging punctuation) and ADD the superscript-figure
% codepoints on top.  Without load=, a global \SetProtrusion REPLACES the
% font's whole protrusion vector, silently disabling all default hanging
% punctuation across the edition.
\edef\verseprot@do@protrusion{%
  \noexpand\SetProtrusion
    [ load = default ]%
    { font = */*/*/*/* }%
    {\verseprot@list}%
}%
\verseprot@do@protrusion

% Adapted for LaTeXGen: called in scripture's font hook AFTER selecting
% scripturefont. The sermon main font is intentionally a different family.
\newif\ifverseprot@use
\newcommand{\detectverseprotrusion}{%
  \setbox0=\hbox{{\addfontfeature{VerticalPosition=Superior}0}}%
  \setbox2=\hbox{0}%
  \ifdim\wd0=\wd2
    \verseprot@usefalse
  \else
    \ifverseprotrusion \verseprot@usetrue \else \verseprot@usefalse \fi
  \fi
}

% --- The formatter wired into scripture's verse/format.
\NewDocumentCommand\verseprotrudenum{m}{%
  \ifverseprot@use
    {\addfontfeature{VerticalPosition=Superior}#1}%
  \else
    \textsuperscript{#1}%
  \fi
}

\makeatother
.}%
}{}

% --- Tunable: partial left-protrusion factor (out of 1000 = one glyph width).
%     Edit ONLY this line to retune the hang; the \edef below propagates it
%     into all ten digit entries automatically.
\def\verseprot@factor{350}

% Build the entry list with the factor expanded in, NO spaces anywhere
% (catcode-10 spaces in the stored token list break \MT@pr@split@val --
%  see header note #4b).  Every newline is suppressed by a trailing `%'.
\edef\verseprot@list{%
"2070={\verseprot@factor,0},%
"00B9={\verseprot@factor,0},%
"00B2={\verseprot@factor,0},%
"00B3={\verseprot@factor,0},%
"2074={\verseprot@factor,0},%
"2075={\verseprot@factor,0},%
"2076={\verseprot@factor,0},%
"2077={\verseprot@factor,0},%
"2078={\verseprot@factor,0},%
"2079={\verseprot@factor,0}}

% --- Manual override (an edition may force the feature off):
\newif\ifverseprotrusion  \verseprotrusiontrue

% --- Global (default) microtype protrusion: left-only hang on the SUPERSCRIPT
%     glyphs that the `sups' feature substitutes for the digits 0-9 (see header
%     note #5). Keyed on the substituted Unicode codepoints, NOT on ASCII 0-9.
%     NO [name=...] and NO [context=...]: this MERGES into the protrusion list
%     microtype loads for the VerticalPosition=Superior font instance, which is
%     what actually binds the codepoints (header note #1). Safe globally because
%     these superscript glyphs only ever appear as verse numbers in this doc.
%     \verseprot@list is passed via \edef so the literal text (not the macro
%     name) is baked into the call before microtype reads it (header note #4).
% [load=default]: INHERIT microtype's built-in protrusion (quotes, hyphens,
% dashes, periods -- the hanging punctuation) and ADD the superscript-figure
% codepoints on top.  Without load=, a global \SetProtrusion REPLACES the
% font's whole protrusion vector, silently disabling all default hanging
% punctuation across the edition.
\edef\verseprot@do@protrusion{%
  \noexpand\SetProtrusion
    [ load = default ]%
    { font = */*/*/*/* }%
    {\verseprot@list}%
}%
\verseprot@do@protrusion

% Adapted for LaTeXGen: called in scripture's font hook AFTER selecting
% scripturefont. The sermon main font is intentionally a different family.
\newif\ifverseprot@use
\newcommand{\detectverseprotrusion}{%
  \setbox0=\hbox{{\addfontfeature{VerticalPosition=Superior}0}}%
  \setbox2=\hbox{0}%
  \ifdim\wd0=\wd2
    \verseprot@usefalse
  \else
    \ifverseprotrusion \verseprot@usetrue \else \verseprot@usefalse \fi
  \fi
}

% --- The formatter wired into scripture's verse/format.
\NewDocumentCommand\verseprotrudenum{m}{%
  \ifverseprot@use
    {\addfontfeature{VerticalPosition=Superior}#1}%
  \else
    \textsuperscript{#1}%
  \fi
}

\makeatother
.}%
}{}
\@ifundefined{addfontfeature}{%
  \PackageError{verse_protrusion}%
    {fontspec must be loaded before verse_protrusion.tex}%
    {Add \string\usepackage{fontspec} (and set a main font) before \string%% verse_protrusion.tex
%% Microtype protrusion for verse numbers (edition-portable).
%% Requires, loaded BEFORE this file: fontspec (main font set), microtype, scripture.
%% Effect: a verse number at the start of a line hangs a tuned partial amount
%% into the left margin. Mid-line numbers are unaffected (microtype only
%% protrudes line-edge glyphs). Font-agnostic: uses OpenType superior figures
%% of the CURRENT body font; falls back to \textsuperscript if unavailable.
%%
%% WORKING MECHANISM (the bits that are easy to get wrong):
%%   1. *** NO named context, NO \microtypecontext. *** Use a PLAIN, default
%%      \SetProtrusion (neither [name=...] nor [context=...]). microtype then
%%      MERGES these entries into whatever protrusion list it loads for any
%%      matching font, so the superior-figure instance gets the codepoints
%%      automatically. The earlier [context=verseprot] +
%%      \microtypecontext{protrusion=verseprot} approach was a NO-OP: the
%%      verbose log showed the context switch loaded an auto-generated EMPTY
%%      list named `<file>/<line>' and the named list never bound to the active
%%      superior font, so lpcode stayed 0 (measured: zero movement at 600 dpi).
%%      The global config is safe here because the superscript-codepoint glyphs
%%      (U+00B9 etc.) appear in this document ONLY as verse numbers (body prose
%%      never emits superior figures), so the entries are effectively scoped to
%%      verse numbers already. PROOF the binding now works: with the global
%%      config, `\the\lpcode\font 185' (U+00B9) reports a nonzero value on the
%%      VerticalPosition=Superior instance (was 0/unset before), and 600 dpi
%%      measurement shows the leading verse number's leftmost ink ~9 px left of
%%      where it sits with the feature off.
%%   2. Font set is the full 5-axis wildcard { font = */*/*/*/* } so it matches
%%      BOTH the body font instance AND the separate font instance that fontspec
%%      spins up for VerticalPosition=Superior (e.g. EBGaramond(1)).
%%   3. Each entry is char = {left,right} (ONE brace, comma-separated). Left-only
%%      hang -> {350,0}. Two separate braces {350}{0} is wrong syntax.
%%   4. To retune the factor, edit \verseprot@factor (one place). The \edef
%%      below expands it into all ten entries before passing to \SetProtrusion.
%%      TWO critical constraints on the expansion approach:
%%      (a) microtype stores \SetProtrusion's third arg via \gdef -- if we pass
%%          `{\verseprot@list}' as the arg, the macro TOKEN (not its expansion)
%%          is stored, and microtype's \MT@set@codes later sees \verseprot@list
%%          as a single opaque token, not as the hex-entry list. We therefore use
%%          `\edef\verseprot@do@protrusion{...{\verseprot@list}...}' to EXPAND the
%%          list inline and store the literal tokens in \verseprot@do@protrusion
%%          before \SetProtrusion is ever called.
%%      (b) microtype reads the \SetProtrusion list with \catcode` '=9 (spaces
%%          ignored), but our \edef runs with normal catcodes where space is
%%          catcode 10. If the expanded list contains catcode-10 space tokens,
%%          \MT@pr@split@val (which reads `left,right\relax') sees a leading space
%%          before `{350,0}' and can't find the `,\relax' terminator -> runaway.
%%          The fix: write the list with NO spaces (all newlines suppressed by `%'
%%          and no space around `=' or `{'). microtype's own parser ignores the
%%          lack of spaces; it just needs the correct token sequence.
%%   5. *** THE OTHER KEY FIX ***  microtype (luatex) keys protrusion by the codepoint
%%      of the glyph AS IT APPEARS IN THE NODE LIST, i.e. AFTER GSUB shaping.
%%      VerticalPosition=Superior runs the OpenType `sups' feature, which in
%%      EB Garamond substitutes each ASCII digit for the glyph at its Unicode
%%      SUPERSCRIPT codepoint:
%%          0->U+2070  1->U+00B9  2->U+00B2  3->U+00B3
%%          4..9 -> U+2074..U+2079
%%      Therefore the list MUST be keyed on those superscript codepoints, NOT on
%%      the ASCII digits 0-9. (Keying 0-9 = U+0030..U+0039 was the bug: those
%%      input glyphs never reach the line because sups replaces them, so the
%%      lpcode never bound to the rendered superior glyph -> ZERO protrusion,
%%      verified by 600 dpi pixel measurement and by dumping the substituted
%%      glyph names with fontTools: three -> uni00B3, etc.)
%%   6. microtype protrudes GLYPH nodes, not boxes. The substituted superior
%%      figures are real glyphs and protrude; \textsuperscript (a raised box)
%%      does not -- hence the OpenType path is required for the hang, with
%%      \textsuperscript only as the non-protruding fallback when the font lacks
%%      superior figures.
%% Verified by measurement (tests/verse_protrusion_visual_check.tex, 600 dpi):
%% with this setup the leading verse number's leftmost ink sits measurably LEFT
%% of the column rule, while mid-line numbers stay inline; feature ON vs OFF
%% now differs at the left edge.
%%
%% PORTABILITY NOTE: the superscript-codepoint mapping above is what EB Garamond
%% (and most fonts following the Unicode/AGL convention) emit for `sups'. A font
%% whose superior figures live at private-use or `.sups'-named glyph slots with
%% no Unicode value would need different keys; in that case microtype cannot
%% address them by codepoint and the \textsuperscript fallback (no protrusion)
%% is the safe outcome.
\makeatletter

% --- Precondition guards: raise a clear error if required packages are absent.
%     Without microtype (\SetProtrusion) or fontspec (\addfontfeature) the file
%     would fail later with a cryptic "Undefined control sequence" inside an \edef.
\@ifundefined{SetProtrusion}{%
  \PackageError{verse_protrusion}%
    {microtype must be loaded before verse_protrusion.tex}%
    {Add \string\usepackage{microtype} (and fontspec) before \string%% verse_protrusion.tex
%% Microtype protrusion for verse numbers (edition-portable).
%% Requires, loaded BEFORE this file: fontspec (main font set), microtype, scripture.
%% Effect: a verse number at the start of a line hangs a tuned partial amount
%% into the left margin. Mid-line numbers are unaffected (microtype only
%% protrudes line-edge glyphs). Font-agnostic: uses OpenType superior figures
%% of the CURRENT body font; falls back to \textsuperscript if unavailable.
%%
%% WORKING MECHANISM (the bits that are easy to get wrong):
%%   1. *** NO named context, NO \microtypecontext. *** Use a PLAIN, default
%%      \SetProtrusion (neither [name=...] nor [context=...]). microtype then
%%      MERGES these entries into whatever protrusion list it loads for any
%%      matching font, so the superior-figure instance gets the codepoints
%%      automatically. The earlier [context=verseprot] +
%%      \microtypecontext{protrusion=verseprot} approach was a NO-OP: the
%%      verbose log showed the context switch loaded an auto-generated EMPTY
%%      list named `<file>/<line>' and the named list never bound to the active
%%      superior font, so lpcode stayed 0 (measured: zero movement at 600 dpi).
%%      The global config is safe here because the superscript-codepoint glyphs
%%      (U+00B9 etc.) appear in this document ONLY as verse numbers (body prose
%%      never emits superior figures), so the entries are effectively scoped to
%%      verse numbers already. PROOF the binding now works: with the global
%%      config, `\the\lpcode\font 185' (U+00B9) reports a nonzero value on the
%%      VerticalPosition=Superior instance (was 0/unset before), and 600 dpi
%%      measurement shows the leading verse number's leftmost ink ~9 px left of
%%      where it sits with the feature off.
%%   2. Font set is the full 5-axis wildcard { font = */*/*/*/* } so it matches
%%      BOTH the body font instance AND the separate font instance that fontspec
%%      spins up for VerticalPosition=Superior (e.g. EBGaramond(1)).
%%   3. Each entry is char = {left,right} (ONE brace, comma-separated). Left-only
%%      hang -> {350,0}. Two separate braces {350}{0} is wrong syntax.
%%   4. To retune the factor, edit \verseprot@factor (one place). The \edef
%%      below expands it into all ten entries before passing to \SetProtrusion.
%%      TWO critical constraints on the expansion approach:
%%      (a) microtype stores \SetProtrusion's third arg via \gdef -- if we pass
%%          `{\verseprot@list}' as the arg, the macro TOKEN (not its expansion)
%%          is stored, and microtype's \MT@set@codes later sees \verseprot@list
%%          as a single opaque token, not as the hex-entry list. We therefore use
%%          `\edef\verseprot@do@protrusion{...{\verseprot@list}...}' to EXPAND the
%%          list inline and store the literal tokens in \verseprot@do@protrusion
%%          before \SetProtrusion is ever called.
%%      (b) microtype reads the \SetProtrusion list with \catcode` '=9 (spaces
%%          ignored), but our \edef runs with normal catcodes where space is
%%          catcode 10. If the expanded list contains catcode-10 space tokens,
%%          \MT@pr@split@val (which reads `left,right\relax') sees a leading space
%%          before `{350,0}' and can't find the `,\relax' terminator -> runaway.
%%          The fix: write the list with NO spaces (all newlines suppressed by `%'
%%          and no space around `=' or `{'). microtype's own parser ignores the
%%          lack of spaces; it just needs the correct token sequence.
%%   5. *** THE OTHER KEY FIX ***  microtype (luatex) keys protrusion by the codepoint
%%      of the glyph AS IT APPEARS IN THE NODE LIST, i.e. AFTER GSUB shaping.
%%      VerticalPosition=Superior runs the OpenType `sups' feature, which in
%%      EB Garamond substitutes each ASCII digit for the glyph at its Unicode
%%      SUPERSCRIPT codepoint:
%%          0->U+2070  1->U+00B9  2->U+00B2  3->U+00B3
%%          4..9 -> U+2074..U+2079
%%      Therefore the list MUST be keyed on those superscript codepoints, NOT on
%%      the ASCII digits 0-9. (Keying 0-9 = U+0030..U+0039 was the bug: those
%%      input glyphs never reach the line because sups replaces them, so the
%%      lpcode never bound to the rendered superior glyph -> ZERO protrusion,
%%      verified by 600 dpi pixel measurement and by dumping the substituted
%%      glyph names with fontTools: three -> uni00B3, etc.)
%%   6. microtype protrudes GLYPH nodes, not boxes. The substituted superior
%%      figures are real glyphs and protrude; \textsuperscript (a raised box)
%%      does not -- hence the OpenType path is required for the hang, with
%%      \textsuperscript only as the non-protruding fallback when the font lacks
%%      superior figures.
%% Verified by measurement (tests/verse_protrusion_visual_check.tex, 600 dpi):
%% with this setup the leading verse number's leftmost ink sits measurably LEFT
%% of the column rule, while mid-line numbers stay inline; feature ON vs OFF
%% now differs at the left edge.
%%
%% PORTABILITY NOTE: the superscript-codepoint mapping above is what EB Garamond
%% (and most fonts following the Unicode/AGL convention) emit for `sups'. A font
%% whose superior figures live at private-use or `.sups'-named glyph slots with
%% no Unicode value would need different keys; in that case microtype cannot
%% address them by codepoint and the \textsuperscript fallback (no protrusion)
%% is the safe outcome.
\makeatletter

% --- Precondition guards: raise a clear error if required packages are absent.
%     Without microtype (\SetProtrusion) or fontspec (\addfontfeature) the file
%     would fail later with a cryptic "Undefined control sequence" inside an \edef.
\@ifundefined{SetProtrusion}{%
  \PackageError{verse_protrusion}%
    {microtype must be loaded before verse_protrusion.tex}%
    {Add \string\usepackage{microtype} (and fontspec) before \string\input{verse_protrusion}.}%
}{}
\@ifundefined{addfontfeature}{%
  \PackageError{verse_protrusion}%
    {fontspec must be loaded before verse_protrusion.tex}%
    {Add \string\usepackage{fontspec} (and set a main font) before \string\input{verse_protrusion}.}%
}{}

% --- Tunable: partial left-protrusion factor (out of 1000 = one glyph width).
%     Edit ONLY this line to retune the hang; the \edef below propagates it
%     into all ten digit entries automatically.
\def\verseprot@factor{350}

% Build the entry list with the factor expanded in, NO spaces anywhere
% (catcode-10 spaces in the stored token list break \MT@pr@split@val --
%  see header note #4b).  Every newline is suppressed by a trailing `%'.
\edef\verseprot@list{%
"2070={\verseprot@factor,0},%
"00B9={\verseprot@factor,0},%
"00B2={\verseprot@factor,0},%
"00B3={\verseprot@factor,0},%
"2074={\verseprot@factor,0},%
"2075={\verseprot@factor,0},%
"2076={\verseprot@factor,0},%
"2077={\verseprot@factor,0},%
"2078={\verseprot@factor,0},%
"2079={\verseprot@factor,0}}

% --- Manual override (an edition may force the feature off):
\newif\ifverseprotrusion  \verseprotrusiontrue

% --- Global (default) microtype protrusion: left-only hang on the SUPERSCRIPT
%     glyphs that the `sups' feature substitutes for the digits 0-9 (see header
%     note #5). Keyed on the substituted Unicode codepoints, NOT on ASCII 0-9.
%     NO [name=...] and NO [context=...]: this MERGES into the protrusion list
%     microtype loads for the VerticalPosition=Superior font instance, which is
%     what actually binds the codepoints (header note #1). Safe globally because
%     these superscript glyphs only ever appear as verse numbers in this doc.
%     \verseprot@list is passed via \edef so the literal text (not the macro
%     name) is baked into the call before microtype reads it (header note #4).
% [load=default]: INHERIT microtype's built-in protrusion (quotes, hyphens,
% dashes, periods -- the hanging punctuation) and ADD the superscript-figure
% codepoints on top.  Without load=, a global \SetProtrusion REPLACES the
% font's whole protrusion vector, silently disabling all default hanging
% punctuation across the edition.
\edef\verseprot@do@protrusion{%
  \noexpand\SetProtrusion
    [ load = default ]%
    { font = */*/*/*/* }%
    {\verseprot@list}%
}%
\verseprot@do@protrusion

% Adapted for LaTeXGen: called in scripture's font hook AFTER selecting
% scripturefont. The sermon main font is intentionally a different family.
\newif\ifverseprot@use
\newcommand{\detectverseprotrusion}{%
  \setbox0=\hbox{{\addfontfeature{VerticalPosition=Superior}0}}%
  \setbox2=\hbox{0}%
  \ifdim\wd0=\wd2
    \verseprot@usefalse
  \else
    \ifverseprotrusion \verseprot@usetrue \else \verseprot@usefalse \fi
  \fi
}

% --- The formatter wired into scripture's verse/format.
\NewDocumentCommand\verseprotrudenum{m}{%
  \ifverseprot@use
    {\addfontfeature{VerticalPosition=Superior}#1}%
  \else
    \textsuperscript{#1}%
  \fi
}

\makeatother
.}%
}{}
\@ifundefined{addfontfeature}{%
  \PackageError{verse_protrusion}%
    {fontspec must be loaded before verse_protrusion.tex}%
    {Add \string\usepackage{fontspec} (and set a main font) before \string%% verse_protrusion.tex
%% Microtype protrusion for verse numbers (edition-portable).
%% Requires, loaded BEFORE this file: fontspec (main font set), microtype, scripture.
%% Effect: a verse number at the start of a line hangs a tuned partial amount
%% into the left margin. Mid-line numbers are unaffected (microtype only
%% protrudes line-edge glyphs). Font-agnostic: uses OpenType superior figures
%% of the CURRENT body font; falls back to \textsuperscript if unavailable.
%%
%% WORKING MECHANISM (the bits that are easy to get wrong):
%%   1. *** NO named context, NO \microtypecontext. *** Use a PLAIN, default
%%      \SetProtrusion (neither [name=...] nor [context=...]). microtype then
%%      MERGES these entries into whatever protrusion list it loads for any
%%      matching font, so the superior-figure instance gets the codepoints
%%      automatically. The earlier [context=verseprot] +
%%      \microtypecontext{protrusion=verseprot} approach was a NO-OP: the
%%      verbose log showed the context switch loaded an auto-generated EMPTY
%%      list named `<file>/<line>' and the named list never bound to the active
%%      superior font, so lpcode stayed 0 (measured: zero movement at 600 dpi).
%%      The global config is safe here because the superscript-codepoint glyphs
%%      (U+00B9 etc.) appear in this document ONLY as verse numbers (body prose
%%      never emits superior figures), so the entries are effectively scoped to
%%      verse numbers already. PROOF the binding now works: with the global
%%      config, `\the\lpcode\font 185' (U+00B9) reports a nonzero value on the
%%      VerticalPosition=Superior instance (was 0/unset before), and 600 dpi
%%      measurement shows the leading verse number's leftmost ink ~9 px left of
%%      where it sits with the feature off.
%%   2. Font set is the full 5-axis wildcard { font = */*/*/*/* } so it matches
%%      BOTH the body font instance AND the separate font instance that fontspec
%%      spins up for VerticalPosition=Superior (e.g. EBGaramond(1)).
%%   3. Each entry is char = {left,right} (ONE brace, comma-separated). Left-only
%%      hang -> {350,0}. Two separate braces {350}{0} is wrong syntax.
%%   4. To retune the factor, edit \verseprot@factor (one place). The \edef
%%      below expands it into all ten entries before passing to \SetProtrusion.
%%      TWO critical constraints on the expansion approach:
%%      (a) microtype stores \SetProtrusion's third arg via \gdef -- if we pass
%%          `{\verseprot@list}' as the arg, the macro TOKEN (not its expansion)
%%          is stored, and microtype's \MT@set@codes later sees \verseprot@list
%%          as a single opaque token, not as the hex-entry list. We therefore use
%%          `\edef\verseprot@do@protrusion{...{\verseprot@list}...}' to EXPAND the
%%          list inline and store the literal tokens in \verseprot@do@protrusion
%%          before \SetProtrusion is ever called.
%%      (b) microtype reads the \SetProtrusion list with \catcode` '=9 (spaces
%%          ignored), but our \edef runs with normal catcodes where space is
%%          catcode 10. If the expanded list contains catcode-10 space tokens,
%%          \MT@pr@split@val (which reads `left,right\relax') sees a leading space
%%          before `{350,0}' and can't find the `,\relax' terminator -> runaway.
%%          The fix: write the list with NO spaces (all newlines suppressed by `%'
%%          and no space around `=' or `{'). microtype's own parser ignores the
%%          lack of spaces; it just needs the correct token sequence.
%%   5. *** THE OTHER KEY FIX ***  microtype (luatex) keys protrusion by the codepoint
%%      of the glyph AS IT APPEARS IN THE NODE LIST, i.e. AFTER GSUB shaping.
%%      VerticalPosition=Superior runs the OpenType `sups' feature, which in
%%      EB Garamond substitutes each ASCII digit for the glyph at its Unicode
%%      SUPERSCRIPT codepoint:
%%          0->U+2070  1->U+00B9  2->U+00B2  3->U+00B3
%%          4..9 -> U+2074..U+2079
%%      Therefore the list MUST be keyed on those superscript codepoints, NOT on
%%      the ASCII digits 0-9. (Keying 0-9 = U+0030..U+0039 was the bug: those
%%      input glyphs never reach the line because sups replaces them, so the
%%      lpcode never bound to the rendered superior glyph -> ZERO protrusion,
%%      verified by 600 dpi pixel measurement and by dumping the substituted
%%      glyph names with fontTools: three -> uni00B3, etc.)
%%   6. microtype protrudes GLYPH nodes, not boxes. The substituted superior
%%      figures are real glyphs and protrude; \textsuperscript (a raised box)
%%      does not -- hence the OpenType path is required for the hang, with
%%      \textsuperscript only as the non-protruding fallback when the font lacks
%%      superior figures.
%% Verified by measurement (tests/verse_protrusion_visual_check.tex, 600 dpi):
%% with this setup the leading verse number's leftmost ink sits measurably LEFT
%% of the column rule, while mid-line numbers stay inline; feature ON vs OFF
%% now differs at the left edge.
%%
%% PORTABILITY NOTE: the superscript-codepoint mapping above is what EB Garamond
%% (and most fonts following the Unicode/AGL convention) emit for `sups'. A font
%% whose superior figures live at private-use or `.sups'-named glyph slots with
%% no Unicode value would need different keys; in that case microtype cannot
%% address them by codepoint and the \textsuperscript fallback (no protrusion)
%% is the safe outcome.
\makeatletter

% --- Precondition guards: raise a clear error if required packages are absent.
%     Without microtype (\SetProtrusion) or fontspec (\addfontfeature) the file
%     would fail later with a cryptic "Undefined control sequence" inside an \edef.
\@ifundefined{SetProtrusion}{%
  \PackageError{verse_protrusion}%
    {microtype must be loaded before verse_protrusion.tex}%
    {Add \string\usepackage{microtype} (and fontspec) before \string\input{verse_protrusion}.}%
}{}
\@ifundefined{addfontfeature}{%
  \PackageError{verse_protrusion}%
    {fontspec must be loaded before verse_protrusion.tex}%
    {Add \string\usepackage{fontspec} (and set a main font) before \string\input{verse_protrusion}.}%
}{}

% --- Tunable: partial left-protrusion factor (out of 1000 = one glyph width).
%     Edit ONLY this line to retune the hang; the \edef below propagates it
%     into all ten digit entries automatically.
\def\verseprot@factor{350}

% Build the entry list with the factor expanded in, NO spaces anywhere
% (catcode-10 spaces in the stored token list break \MT@pr@split@val --
%  see header note #4b).  Every newline is suppressed by a trailing `%'.
\edef\verseprot@list{%
"2070={\verseprot@factor,0},%
"00B9={\verseprot@factor,0},%
"00B2={\verseprot@factor,0},%
"00B3={\verseprot@factor,0},%
"2074={\verseprot@factor,0},%
"2075={\verseprot@factor,0},%
"2076={\verseprot@factor,0},%
"2077={\verseprot@factor,0},%
"2078={\verseprot@factor,0},%
"2079={\verseprot@factor,0}}

% --- Manual override (an edition may force the feature off):
\newif\ifverseprotrusion  \verseprotrusiontrue

% --- Global (default) microtype protrusion: left-only hang on the SUPERSCRIPT
%     glyphs that the `sups' feature substitutes for the digits 0-9 (see header
%     note #5). Keyed on the substituted Unicode codepoints, NOT on ASCII 0-9.
%     NO [name=...] and NO [context=...]: this MERGES into the protrusion list
%     microtype loads for the VerticalPosition=Superior font instance, which is
%     what actually binds the codepoints (header note #1). Safe globally because
%     these superscript glyphs only ever appear as verse numbers in this doc.
%     \verseprot@list is passed via \edef so the literal text (not the macro
%     name) is baked into the call before microtype reads it (header note #4).
% [load=default]: INHERIT microtype's built-in protrusion (quotes, hyphens,
% dashes, periods -- the hanging punctuation) and ADD the superscript-figure
% codepoints on top.  Without load=, a global \SetProtrusion REPLACES the
% font's whole protrusion vector, silently disabling all default hanging
% punctuation across the edition.
\edef\verseprot@do@protrusion{%
  \noexpand\SetProtrusion
    [ load = default ]%
    { font = */*/*/*/* }%
    {\verseprot@list}%
}%
\verseprot@do@protrusion

% Adapted for LaTeXGen: called in scripture's font hook AFTER selecting
% scripturefont. The sermon main font is intentionally a different family.
\newif\ifverseprot@use
\newcommand{\detectverseprotrusion}{%
  \setbox0=\hbox{{\addfontfeature{VerticalPosition=Superior}0}}%
  \setbox2=\hbox{0}%
  \ifdim\wd0=\wd2
    \verseprot@usefalse
  \else
    \ifverseprotrusion \verseprot@usetrue \else \verseprot@usefalse \fi
  \fi
}

% --- The formatter wired into scripture's verse/format.
\NewDocumentCommand\verseprotrudenum{m}{%
  \ifverseprot@use
    {\addfontfeature{VerticalPosition=Superior}#1}%
  \else
    \textsuperscript{#1}%
  \fi
}

\makeatother
.}%
}{}

% --- Tunable: partial left-protrusion factor (out of 1000 = one glyph width).
%     Edit ONLY this line to retune the hang; the \edef below propagates it
%     into all ten digit entries automatically.
\def\verseprot@factor{350}

% Build the entry list with the factor expanded in, NO spaces anywhere
% (catcode-10 spaces in the stored token list break \MT@pr@split@val --
%  see header note #4b).  Every newline is suppressed by a trailing `%'.
\edef\verseprot@list{%
"2070={\verseprot@factor,0},%
"00B9={\verseprot@factor,0},%
"00B2={\verseprot@factor,0},%
"00B3={\verseprot@factor,0},%
"2074={\verseprot@factor,0},%
"2075={\verseprot@factor,0},%
"2076={\verseprot@factor,0},%
"2077={\verseprot@factor,0},%
"2078={\verseprot@factor,0},%
"2079={\verseprot@factor,0}}

% --- Manual override (an edition may force the feature off):
\newif\ifverseprotrusion  \verseprotrusiontrue

% --- Global (default) microtype protrusion: left-only hang on the SUPERSCRIPT
%     glyphs that the `sups' feature substitutes for the digits 0-9 (see header
%     note #5). Keyed on the substituted Unicode codepoints, NOT on ASCII 0-9.
%     NO [name=...] and NO [context=...]: this MERGES into the protrusion list
%     microtype loads for the VerticalPosition=Superior font instance, which is
%     what actually binds the codepoints (header note #1). Safe globally because
%     these superscript glyphs only ever appear as verse numbers in this doc.
%     \verseprot@list is passed via \edef so the literal text (not the macro
%     name) is baked into the call before microtype reads it (header note #4).
% [load=default]: INHERIT microtype's built-in protrusion (quotes, hyphens,
% dashes, periods -- the hanging punctuation) and ADD the superscript-figure
% codepoints on top.  Without load=, a global \SetProtrusion REPLACES the
% font's whole protrusion vector, silently disabling all default hanging
% punctuation across the edition.
\edef\verseprot@do@protrusion{%
  \noexpand\SetProtrusion
    [ load = default ]%
    { font = */*/*/*/* }%
    {\verseprot@list}%
}%
\verseprot@do@protrusion

% Adapted for LaTeXGen: called in scripture's font hook AFTER selecting
% scripturefont. The sermon main font is intentionally a different family.
\newif\ifverseprot@use
\newcommand{\detectverseprotrusion}{%
  \setbox0=\hbox{{\addfontfeature{VerticalPosition=Superior}0}}%
  \setbox2=\hbox{0}%
  \ifdim\wd0=\wd2
    \verseprot@usefalse
  \else
    \ifverseprotrusion \verseprot@usetrue \else \verseprot@usefalse \fi
  \fi
}

% --- The formatter wired into scripture's verse/format.
\NewDocumentCommand\verseprotrudenum{m}{%
  \ifverseprot@use
    {\addfontfeature{VerticalPosition=Superior}#1}%
  \else
    \textsuperscript{#1}%
  \fi
}

\makeatother
.}%
}{}

% --- Tunable: partial left-protrusion factor (out of 1000 = one glyph width).
%     Edit ONLY this line to retune the hang; the \edef below propagates it
%     into all ten digit entries automatically.
\def\verseprot@factor{350}

% Build the entry list with the factor expanded in, NO spaces anywhere
% (catcode-10 spaces in the stored token list break \MT@pr@split@val --
%  see header note #4b).  Every newline is suppressed by a trailing `%'.
\edef\verseprot@list{%
"2070={\verseprot@factor,0},%
"00B9={\verseprot@factor,0},%
"00B2={\verseprot@factor,0},%
"00B3={\verseprot@factor,0},%
"2074={\verseprot@factor,0},%
"2075={\verseprot@factor,0},%
"2076={\verseprot@factor,0},%
"2077={\verseprot@factor,0},%
"2078={\verseprot@factor,0},%
"2079={\verseprot@factor,0}}

% --- Manual override (an edition may force the feature off):
\newif\ifverseprotrusion  \verseprotrusiontrue

% --- Global (default) microtype protrusion: left-only hang on the SUPERSCRIPT
%     glyphs that the `sups' feature substitutes for the digits 0-9 (see header
%     note #5). Keyed on the substituted Unicode codepoints, NOT on ASCII 0-9.
%     NO [name=...] and NO [context=...]: this MERGES into the protrusion list
%     microtype loads for the VerticalPosition=Superior font instance, which is
%     what actually binds the codepoints (header note #1). Safe globally because
%     these superscript glyphs only ever appear as verse numbers in this doc.
%     \verseprot@list is passed via \edef so the literal text (not the macro
%     name) is baked into the call before microtype reads it (header note #4).
% [load=default]: INHERIT microtype's built-in protrusion (quotes, hyphens,
% dashes, periods -- the hanging punctuation) and ADD the superscript-figure
% codepoints on top.  Without load=, a global \SetProtrusion REPLACES the
% font's whole protrusion vector, silently disabling all default hanging
% punctuation across the edition.
\edef\verseprot@do@protrusion{%
  \noexpand\SetProtrusion
    [ load = default ]%
    { font = */*/*/*/* }%
    {\verseprot@list}%
}%
\verseprot@do@protrusion

% Adapted for LaTeXGen: called in scripture's font hook AFTER selecting
% scripturefont. The sermon main font is intentionally a different family.
\newif\ifverseprot@use
\newcommand{\detectverseprotrusion}{%
  \setbox0=\hbox{{\addfontfeature{VerticalPosition=Superior}0}}%
  \setbox2=\hbox{0}%
  \ifdim\wd0=\wd2
    \verseprot@usefalse
  \else
    \ifverseprotrusion \verseprot@usetrue \else \verseprot@usefalse \fi
  \fi
}

% --- The formatter wired into scripture's verse/format.
\NewDocumentCommand\verseprotrudenum{m}{%
  \ifverseprot@use
    {\addfontfeature{VerticalPosition=Superior}#1}%
  \else
    \textsuperscript{#1}%
  \fi
}

\makeatother
.}%
}{}
\@ifundefined{addfontfeature}{%
  \PackageError{verse_protrusion}%
    {fontspec must be loaded before verse_protrusion.tex}%
    {Add \string\usepackage{fontspec} (and set a main font) before \string%% verse_protrusion.tex
%% Microtype protrusion for verse numbers (edition-portable).
%% Requires, loaded BEFORE this file: fontspec (main font set), microtype, scripture.
%% Effect: a verse number at the start of a line hangs a tuned partial amount
%% into the left margin. Mid-line numbers are unaffected (microtype only
%% protrudes line-edge glyphs). Font-agnostic: uses OpenType superior figures
%% of the CURRENT body font; falls back to \textsuperscript if unavailable.
%%
%% WORKING MECHANISM (the bits that are easy to get wrong):
%%   1. *** NO named context, NO \microtypecontext. *** Use a PLAIN, default
%%      \SetProtrusion (neither [name=...] nor [context=...]). microtype then
%%      MERGES these entries into whatever protrusion list it loads for any
%%      matching font, so the superior-figure instance gets the codepoints
%%      automatically. The earlier [context=verseprot] +
%%      \microtypecontext{protrusion=verseprot} approach was a NO-OP: the
%%      verbose log showed the context switch loaded an auto-generated EMPTY
%%      list named `<file>/<line>' and the named list never bound to the active
%%      superior font, so lpcode stayed 0 (measured: zero movement at 600 dpi).
%%      The global config is safe here because the superscript-codepoint glyphs
%%      (U+00B9 etc.) appear in this document ONLY as verse numbers (body prose
%%      never emits superior figures), so the entries are effectively scoped to
%%      verse numbers already. PROOF the binding now works: with the global
%%      config, `\the\lpcode\font 185' (U+00B9) reports a nonzero value on the
%%      VerticalPosition=Superior instance (was 0/unset before), and 600 dpi
%%      measurement shows the leading verse number's leftmost ink ~9 px left of
%%      where it sits with the feature off.
%%   2. Font set is the full 5-axis wildcard { font = */*/*/*/* } so it matches
%%      BOTH the body font instance AND the separate font instance that fontspec
%%      spins up for VerticalPosition=Superior (e.g. EBGaramond(1)).
%%   3. Each entry is char = {left,right} (ONE brace, comma-separated). Left-only
%%      hang -> {350,0}. Two separate braces {350}{0} is wrong syntax.
%%   4. To retune the factor, edit \verseprot@factor (one place). The \edef
%%      below expands it into all ten entries before passing to \SetProtrusion.
%%      TWO critical constraints on the expansion approach:
%%      (a) microtype stores \SetProtrusion's third arg via \gdef -- if we pass
%%          `{\verseprot@list}' as the arg, the macro TOKEN (not its expansion)
%%          is stored, and microtype's \MT@set@codes later sees \verseprot@list
%%          as a single opaque token, not as the hex-entry list. We therefore use
%%          `\edef\verseprot@do@protrusion{...{\verseprot@list}...}' to EXPAND the
%%          list inline and store the literal tokens in \verseprot@do@protrusion
%%          before \SetProtrusion is ever called.
%%      (b) microtype reads the \SetProtrusion list with \catcode` '=9 (spaces
%%          ignored), but our \edef runs with normal catcodes where space is
%%          catcode 10. If the expanded list contains catcode-10 space tokens,
%%          \MT@pr@split@val (which reads `left,right\relax') sees a leading space
%%          before `{350,0}' and can't find the `,\relax' terminator -> runaway.
%%          The fix: write the list with NO spaces (all newlines suppressed by `%'
%%          and no space around `=' or `{'). microtype's own parser ignores the
%%          lack of spaces; it just needs the correct token sequence.
%%   5. *** THE OTHER KEY FIX ***  microtype (luatex) keys protrusion by the codepoint
%%      of the glyph AS IT APPEARS IN THE NODE LIST, i.e. AFTER GSUB shaping.
%%      VerticalPosition=Superior runs the OpenType `sups' feature, which in
%%      EB Garamond substitutes each ASCII digit for the glyph at its Unicode
%%      SUPERSCRIPT codepoint:
%%          0->U+2070  1->U+00B9  2->U+00B2  3->U+00B3
%%          4..9 -> U+2074..U+2079
%%      Therefore the list MUST be keyed on those superscript codepoints, NOT on
%%      the ASCII digits 0-9. (Keying 0-9 = U+0030..U+0039 was the bug: those
%%      input glyphs never reach the line because sups replaces them, so the
%%      lpcode never bound to the rendered superior glyph -> ZERO protrusion,
%%      verified by 600 dpi pixel measurement and by dumping the substituted
%%      glyph names with fontTools: three -> uni00B3, etc.)
%%   6. microtype protrudes GLYPH nodes, not boxes. The substituted superior
%%      figures are real glyphs and protrude; \textsuperscript (a raised box)
%%      does not -- hence the OpenType path is required for the hang, with
%%      \textsuperscript only as the non-protruding fallback when the font lacks
%%      superior figures.
%% Verified by measurement (tests/verse_protrusion_visual_check.tex, 600 dpi):
%% with this setup the leading verse number's leftmost ink sits measurably LEFT
%% of the column rule, while mid-line numbers stay inline; feature ON vs OFF
%% now differs at the left edge.
%%
%% PORTABILITY NOTE: the superscript-codepoint mapping above is what EB Garamond
%% (and most fonts following the Unicode/AGL convention) emit for `sups'. A font
%% whose superior figures live at private-use or `.sups'-named glyph slots with
%% no Unicode value would need different keys; in that case microtype cannot
%% address them by codepoint and the \textsuperscript fallback (no protrusion)
%% is the safe outcome.
\makeatletter

% --- Precondition guards: raise a clear error if required packages are absent.
%     Without microtype (\SetProtrusion) or fontspec (\addfontfeature) the file
%     would fail later with a cryptic "Undefined control sequence" inside an \edef.
\@ifundefined{SetProtrusion}{%
  \PackageError{verse_protrusion}%
    {microtype must be loaded before verse_protrusion.tex}%
    {Add \string\usepackage{microtype} (and fontspec) before \string%% verse_protrusion.tex
%% Microtype protrusion for verse numbers (edition-portable).
%% Requires, loaded BEFORE this file: fontspec (main font set), microtype, scripture.
%% Effect: a verse number at the start of a line hangs a tuned partial amount
%% into the left margin. Mid-line numbers are unaffected (microtype only
%% protrudes line-edge glyphs). Font-agnostic: uses OpenType superior figures
%% of the CURRENT body font; falls back to \textsuperscript if unavailable.
%%
%% WORKING MECHANISM (the bits that are easy to get wrong):
%%   1. *** NO named context, NO \microtypecontext. *** Use a PLAIN, default
%%      \SetProtrusion (neither [name=...] nor [context=...]). microtype then
%%      MERGES these entries into whatever protrusion list it loads for any
%%      matching font, so the superior-figure instance gets the codepoints
%%      automatically. The earlier [context=verseprot] +
%%      \microtypecontext{protrusion=verseprot} approach was a NO-OP: the
%%      verbose log showed the context switch loaded an auto-generated EMPTY
%%      list named `<file>/<line>' and the named list never bound to the active
%%      superior font, so lpcode stayed 0 (measured: zero movement at 600 dpi).
%%      The global config is safe here because the superscript-codepoint glyphs
%%      (U+00B9 etc.) appear in this document ONLY as verse numbers (body prose
%%      never emits superior figures), so the entries are effectively scoped to
%%      verse numbers already. PROOF the binding now works: with the global
%%      config, `\the\lpcode\font 185' (U+00B9) reports a nonzero value on the
%%      VerticalPosition=Superior instance (was 0/unset before), and 600 dpi
%%      measurement shows the leading verse number's leftmost ink ~9 px left of
%%      where it sits with the feature off.
%%   2. Font set is the full 5-axis wildcard { font = */*/*/*/* } so it matches
%%      BOTH the body font instance AND the separate font instance that fontspec
%%      spins up for VerticalPosition=Superior (e.g. EBGaramond(1)).
%%   3. Each entry is char = {left,right} (ONE brace, comma-separated). Left-only
%%      hang -> {350,0}. Two separate braces {350}{0} is wrong syntax.
%%   4. To retune the factor, edit \verseprot@factor (one place). The \edef
%%      below expands it into all ten entries before passing to \SetProtrusion.
%%      TWO critical constraints on the expansion approach:
%%      (a) microtype stores \SetProtrusion's third arg via \gdef -- if we pass
%%          `{\verseprot@list}' as the arg, the macro TOKEN (not its expansion)
%%          is stored, and microtype's \MT@set@codes later sees \verseprot@list
%%          as a single opaque token, not as the hex-entry list. We therefore use
%%          `\edef\verseprot@do@protrusion{...{\verseprot@list}...}' to EXPAND the
%%          list inline and store the literal tokens in \verseprot@do@protrusion
%%          before \SetProtrusion is ever called.
%%      (b) microtype reads the \SetProtrusion list with \catcode` '=9 (spaces
%%          ignored), but our \edef runs with normal catcodes where space is
%%          catcode 10. If the expanded list contains catcode-10 space tokens,
%%          \MT@pr@split@val (which reads `left,right\relax') sees a leading space
%%          before `{350,0}' and can't find the `,\relax' terminator -> runaway.
%%          The fix: write the list with NO spaces (all newlines suppressed by `%'
%%          and no space around `=' or `{'). microtype's own parser ignores the
%%          lack of spaces; it just needs the correct token sequence.
%%   5. *** THE OTHER KEY FIX ***  microtype (luatex) keys protrusion by the codepoint
%%      of the glyph AS IT APPEARS IN THE NODE LIST, i.e. AFTER GSUB shaping.
%%      VerticalPosition=Superior runs the OpenType `sups' feature, which in
%%      EB Garamond substitutes each ASCII digit for the glyph at its Unicode
%%      SUPERSCRIPT codepoint:
%%          0->U+2070  1->U+00B9  2->U+00B2  3->U+00B3
%%          4..9 -> U+2074..U+2079
%%      Therefore the list MUST be keyed on those superscript codepoints, NOT on
%%      the ASCII digits 0-9. (Keying 0-9 = U+0030..U+0039 was the bug: those
%%      input glyphs never reach the line because sups replaces them, so the
%%      lpcode never bound to the rendered superior glyph -> ZERO protrusion,
%%      verified by 600 dpi pixel measurement and by dumping the substituted
%%      glyph names with fontTools: three -> uni00B3, etc.)
%%   6. microtype protrudes GLYPH nodes, not boxes. The substituted superior
%%      figures are real glyphs and protrude; \textsuperscript (a raised box)
%%      does not -- hence the OpenType path is required for the hang, with
%%      \textsuperscript only as the non-protruding fallback when the font lacks
%%      superior figures.
%% Verified by measurement (tests/verse_protrusion_visual_check.tex, 600 dpi):
%% with this setup the leading verse number's leftmost ink sits measurably LEFT
%% of the column rule, while mid-line numbers stay inline; feature ON vs OFF
%% now differs at the left edge.
%%
%% PORTABILITY NOTE: the superscript-codepoint mapping above is what EB Garamond
%% (and most fonts following the Unicode/AGL convention) emit for `sups'. A font
%% whose superior figures live at private-use or `.sups'-named glyph slots with
%% no Unicode value would need different keys; in that case microtype cannot
%% address them by codepoint and the \textsuperscript fallback (no protrusion)
%% is the safe outcome.
\makeatletter

% --- Precondition guards: raise a clear error if required packages are absent.
%     Without microtype (\SetProtrusion) or fontspec (\addfontfeature) the file
%     would fail later with a cryptic "Undefined control sequence" inside an \edef.
\@ifundefined{SetProtrusion}{%
  \PackageError{verse_protrusion}%
    {microtype must be loaded before verse_protrusion.tex}%
    {Add \string\usepackage{microtype} (and fontspec) before \string%% verse_protrusion.tex
%% Microtype protrusion for verse numbers (edition-portable).
%% Requires, loaded BEFORE this file: fontspec (main font set), microtype, scripture.
%% Effect: a verse number at the start of a line hangs a tuned partial amount
%% into the left margin. Mid-line numbers are unaffected (microtype only
%% protrudes line-edge glyphs). Font-agnostic: uses OpenType superior figures
%% of the CURRENT body font; falls back to \textsuperscript if unavailable.
%%
%% WORKING MECHANISM (the bits that are easy to get wrong):
%%   1. *** NO named context, NO \microtypecontext. *** Use a PLAIN, default
%%      \SetProtrusion (neither [name=...] nor [context=...]). microtype then
%%      MERGES these entries into whatever protrusion list it loads for any
%%      matching font, so the superior-figure instance gets the codepoints
%%      automatically. The earlier [context=verseprot] +
%%      \microtypecontext{protrusion=verseprot} approach was a NO-OP: the
%%      verbose log showed the context switch loaded an auto-generated EMPTY
%%      list named `<file>/<line>' and the named list never bound to the active
%%      superior font, so lpcode stayed 0 (measured: zero movement at 600 dpi).
%%      The global config is safe here because the superscript-codepoint glyphs
%%      (U+00B9 etc.) appear in this document ONLY as verse numbers (body prose
%%      never emits superior figures), so the entries are effectively scoped to
%%      verse numbers already. PROOF the binding now works: with the global
%%      config, `\the\lpcode\font 185' (U+00B9) reports a nonzero value on the
%%      VerticalPosition=Superior instance (was 0/unset before), and 600 dpi
%%      measurement shows the leading verse number's leftmost ink ~9 px left of
%%      where it sits with the feature off.
%%   2. Font set is the full 5-axis wildcard { font = */*/*/*/* } so it matches
%%      BOTH the body font instance AND the separate font instance that fontspec
%%      spins up for VerticalPosition=Superior (e.g. EBGaramond(1)).
%%   3. Each entry is char = {left,right} (ONE brace, comma-separated). Left-only
%%      hang -> {350,0}. Two separate braces {350}{0} is wrong syntax.
%%   4. To retune the factor, edit \verseprot@factor (one place). The \edef
%%      below expands it into all ten entries before passing to \SetProtrusion.
%%      TWO critical constraints on the expansion approach:
%%      (a) microtype stores \SetProtrusion's third arg via \gdef -- if we pass
%%          `{\verseprot@list}' as the arg, the macro TOKEN (not its expansion)
%%          is stored, and microtype's \MT@set@codes later sees \verseprot@list
%%          as a single opaque token, not as the hex-entry list. We therefore use
%%          `\edef\verseprot@do@protrusion{...{\verseprot@list}...}' to EXPAND the
%%          list inline and store the literal tokens in \verseprot@do@protrusion
%%          before \SetProtrusion is ever called.
%%      (b) microtype reads the \SetProtrusion list with \catcode` '=9 (spaces
%%          ignored), but our \edef runs with normal catcodes where space is
%%          catcode 10. If the expanded list contains catcode-10 space tokens,
%%          \MT@pr@split@val (which reads `left,right\relax') sees a leading space
%%          before `{350,0}' and can't find the `,\relax' terminator -> runaway.
%%          The fix: write the list with NO spaces (all newlines suppressed by `%'
%%          and no space around `=' or `{'). microtype's own parser ignores the
%%          lack of spaces; it just needs the correct token sequence.
%%   5. *** THE OTHER KEY FIX ***  microtype (luatex) keys protrusion by the codepoint
%%      of the glyph AS IT APPEARS IN THE NODE LIST, i.e. AFTER GSUB shaping.
%%      VerticalPosition=Superior runs the OpenType `sups' feature, which in
%%      EB Garamond substitutes each ASCII digit for the glyph at its Unicode
%%      SUPERSCRIPT codepoint:
%%          0->U+2070  1->U+00B9  2->U+00B2  3->U+00B3
%%          4..9 -> U+2074..U+2079
%%      Therefore the list MUST be keyed on those superscript codepoints, NOT on
%%      the ASCII digits 0-9. (Keying 0-9 = U+0030..U+0039 was the bug: those
%%      input glyphs never reach the line because sups replaces them, so the
%%      lpcode never bound to the rendered superior glyph -> ZERO protrusion,
%%      verified by 600 dpi pixel measurement and by dumping the substituted
%%      glyph names with fontTools: three -> uni00B3, etc.)
%%   6. microtype protrudes GLYPH nodes, not boxes. The substituted superior
%%      figures are real glyphs and protrude; \textsuperscript (a raised box)
%%      does not -- hence the OpenType path is required for the hang, with
%%      \textsuperscript only as the non-protruding fallback when the font lacks
%%      superior figures.
%% Verified by measurement (tests/verse_protrusion_visual_check.tex, 600 dpi):
%% with this setup the leading verse number's leftmost ink sits measurably LEFT
%% of the column rule, while mid-line numbers stay inline; feature ON vs OFF
%% now differs at the left edge.
%%
%% PORTABILITY NOTE: the superscript-codepoint mapping above is what EB Garamond
%% (and most fonts following the Unicode/AGL convention) emit for `sups'. A font
%% whose superior figures live at private-use or `.sups'-named glyph slots with
%% no Unicode value would need different keys; in that case microtype cannot
%% address them by codepoint and the \textsuperscript fallback (no protrusion)
%% is the safe outcome.
\makeatletter

% --- Precondition guards: raise a clear error if required packages are absent.
%     Without microtype (\SetProtrusion) or fontspec (\addfontfeature) the file
%     would fail later with a cryptic "Undefined control sequence" inside an \edef.
\@ifundefined{SetProtrusion}{%
  \PackageError{verse_protrusion}%
    {microtype must be loaded before verse_protrusion.tex}%
    {Add \string\usepackage{microtype} (and fontspec) before \string\input{verse_protrusion}.}%
}{}
\@ifundefined{addfontfeature}{%
  \PackageError{verse_protrusion}%
    {fontspec must be loaded before verse_protrusion.tex}%
    {Add \string\usepackage{fontspec} (and set a main font) before \string\input{verse_protrusion}.}%
}{}

% --- Tunable: partial left-protrusion factor (out of 1000 = one glyph width).
%     Edit ONLY this line to retune the hang; the \edef below propagates it
%     into all ten digit entries automatically.
\def\verseprot@factor{350}

% Build the entry list with the factor expanded in, NO spaces anywhere
% (catcode-10 spaces in the stored token list break \MT@pr@split@val --
%  see header note #4b).  Every newline is suppressed by a trailing `%'.
\edef\verseprot@list{%
"2070={\verseprot@factor,0},%
"00B9={\verseprot@factor,0},%
"00B2={\verseprot@factor,0},%
"00B3={\verseprot@factor,0},%
"2074={\verseprot@factor,0},%
"2075={\verseprot@factor,0},%
"2076={\verseprot@factor,0},%
"2077={\verseprot@factor,0},%
"2078={\verseprot@factor,0},%
"2079={\verseprot@factor,0}}

% --- Manual override (an edition may force the feature off):
\newif\ifverseprotrusion  \verseprotrusiontrue

% --- Global (default) microtype protrusion: left-only hang on the SUPERSCRIPT
%     glyphs that the `sups' feature substitutes for the digits 0-9 (see header
%     note #5). Keyed on the substituted Unicode codepoints, NOT on ASCII 0-9.
%     NO [name=...] and NO [context=...]: this MERGES into the protrusion list
%     microtype loads for the VerticalPosition=Superior font instance, which is
%     what actually binds the codepoints (header note #1). Safe globally because
%     these superscript glyphs only ever appear as verse numbers in this doc.
%     \verseprot@list is passed via \edef so the literal text (not the macro
%     name) is baked into the call before microtype reads it (header note #4).
% [load=default]: INHERIT microtype's built-in protrusion (quotes, hyphens,
% dashes, periods -- the hanging punctuation) and ADD the superscript-figure
% codepoints on top.  Without load=, a global \SetProtrusion REPLACES the
% font's whole protrusion vector, silently disabling all default hanging
% punctuation across the edition.
\edef\verseprot@do@protrusion{%
  \noexpand\SetProtrusion
    [ load = default ]%
    { font = */*/*/*/* }%
    {\verseprot@list}%
}%
\verseprot@do@protrusion

% Adapted for LaTeXGen: called in scripture's font hook AFTER selecting
% scripturefont. The sermon main font is intentionally a different family.
\newif\ifverseprot@use
\newcommand{\detectverseprotrusion}{%
  \setbox0=\hbox{{\addfontfeature{VerticalPosition=Superior}0}}%
  \setbox2=\hbox{0}%
  \ifdim\wd0=\wd2
    \verseprot@usefalse
  \else
    \ifverseprotrusion \verseprot@usetrue \else \verseprot@usefalse \fi
  \fi
}

% --- The formatter wired into scripture's verse/format.
\NewDocumentCommand\verseprotrudenum{m}{%
  \ifverseprot@use
    {\addfontfeature{VerticalPosition=Superior}#1}%
  \else
    \textsuperscript{#1}%
  \fi
}

\makeatother
.}%
}{}
\@ifundefined{addfontfeature}{%
  \PackageError{verse_protrusion}%
    {fontspec must be loaded before verse_protrusion.tex}%
    {Add \string\usepackage{fontspec} (and set a main font) before \string%% verse_protrusion.tex
%% Microtype protrusion for verse numbers (edition-portable).
%% Requires, loaded BEFORE this file: fontspec (main font set), microtype, scripture.
%% Effect: a verse number at the start of a line hangs a tuned partial amount
%% into the left margin. Mid-line numbers are unaffected (microtype only
%% protrudes line-edge glyphs). Font-agnostic: uses OpenType superior figures
%% of the CURRENT body font; falls back to \textsuperscript if unavailable.
%%
%% WORKING MECHANISM (the bits that are easy to get wrong):
%%   1. *** NO named context, NO \microtypecontext. *** Use a PLAIN, default
%%      \SetProtrusion (neither [name=...] nor [context=...]). microtype then
%%      MERGES these entries into whatever protrusion list it loads for any
%%      matching font, so the superior-figure instance gets the codepoints
%%      automatically. The earlier [context=verseprot] +
%%      \microtypecontext{protrusion=verseprot} approach was a NO-OP: the
%%      verbose log showed the context switch loaded an auto-generated EMPTY
%%      list named `<file>/<line>' and the named list never bound to the active
%%      superior font, so lpcode stayed 0 (measured: zero movement at 600 dpi).
%%      The global config is safe here because the superscript-codepoint glyphs
%%      (U+00B9 etc.) appear in this document ONLY as verse numbers (body prose
%%      never emits superior figures), so the entries are effectively scoped to
%%      verse numbers already. PROOF the binding now works: with the global
%%      config, `\the\lpcode\font 185' (U+00B9) reports a nonzero value on the
%%      VerticalPosition=Superior instance (was 0/unset before), and 600 dpi
%%      measurement shows the leading verse number's leftmost ink ~9 px left of
%%      where it sits with the feature off.
%%   2. Font set is the full 5-axis wildcard { font = */*/*/*/* } so it matches
%%      BOTH the body font instance AND the separate font instance that fontspec
%%      spins up for VerticalPosition=Superior (e.g. EBGaramond(1)).
%%   3. Each entry is char = {left,right} (ONE brace, comma-separated). Left-only
%%      hang -> {350,0}. Two separate braces {350}{0} is wrong syntax.
%%   4. To retune the factor, edit \verseprot@factor (one place). The \edef
%%      below expands it into all ten entries before passing to \SetProtrusion.
%%      TWO critical constraints on the expansion approach:
%%      (a) microtype stores \SetProtrusion's third arg via \gdef -- if we pass
%%          `{\verseprot@list}' as the arg, the macro TOKEN (not its expansion)
%%          is stored, and microtype's \MT@set@codes later sees \verseprot@list
%%          as a single opaque token, not as the hex-entry list. We therefore use
%%          `\edef\verseprot@do@protrusion{...{\verseprot@list}...}' to EXPAND the
%%          list inline and store the literal tokens in \verseprot@do@protrusion
%%          before \SetProtrusion is ever called.
%%      (b) microtype reads the \SetProtrusion list with \catcode` '=9 (spaces
%%          ignored), but our \edef runs with normal catcodes where space is
%%          catcode 10. If the expanded list contains catcode-10 space tokens,
%%          \MT@pr@split@val (which reads `left,right\relax') sees a leading space
%%          before `{350,0}' and can't find the `,\relax' terminator -> runaway.
%%          The fix: write the list with NO spaces (all newlines suppressed by `%'
%%          and no space around `=' or `{'). microtype's own parser ignores the
%%          lack of spaces; it just needs the correct token sequence.
%%   5. *** THE OTHER KEY FIX ***  microtype (luatex) keys protrusion by the codepoint
%%      of the glyph AS IT APPEARS IN THE NODE LIST, i.e. AFTER GSUB shaping.
%%      VerticalPosition=Superior runs the OpenType `sups' feature, which in
%%      EB Garamond substitutes each ASCII digit for the glyph at its Unicode
%%      SUPERSCRIPT codepoint:
%%          0->U+2070  1->U+00B9  2->U+00B2  3->U+00B3
%%          4..9 -> U+2074..U+2079
%%      Therefore the list MUST be keyed on those superscript codepoints, NOT on
%%      the ASCII digits 0-9. (Keying 0-9 = U+0030..U+0039 was the bug: those
%%      input glyphs never reach the line because sups replaces them, so the
%%      lpcode never bound to the rendered superior glyph -> ZERO protrusion,
%%      verified by 600 dpi pixel measurement and by dumping the substituted
%%      glyph names with fontTools: three -> uni00B3, etc.)
%%   6. microtype protrudes GLYPH nodes, not boxes. The substituted superior
%%      figures are real glyphs and protrude; \textsuperscript (a raised box)
%%      does not -- hence the OpenType path is required for the hang, with
%%      \textsuperscript only as the non-protruding fallback when the font lacks
%%      superior figures.
%% Verified by measurement (tests/verse_protrusion_visual_check.tex, 600 dpi):
%% with this setup the leading verse number's leftmost ink sits measurably LEFT
%% of the column rule, while mid-line numbers stay inline; feature ON vs OFF
%% now differs at the left edge.
%%
%% PORTABILITY NOTE: the superscript-codepoint mapping above is what EB Garamond
%% (and most fonts following the Unicode/AGL convention) emit for `sups'. A font
%% whose superior figures live at private-use or `.sups'-named glyph slots with
%% no Unicode value would need different keys; in that case microtype cannot
%% address them by codepoint and the \textsuperscript fallback (no protrusion)
%% is the safe outcome.
\makeatletter

% --- Precondition guards: raise a clear error if required packages are absent.
%     Without microtype (\SetProtrusion) or fontspec (\addfontfeature) the file
%     would fail later with a cryptic "Undefined control sequence" inside an \edef.
\@ifundefined{SetProtrusion}{%
  \PackageError{verse_protrusion}%
    {microtype must be loaded before verse_protrusion.tex}%
    {Add \string\usepackage{microtype} (and fontspec) before \string\input{verse_protrusion}.}%
}{}
\@ifundefined{addfontfeature}{%
  \PackageError{verse_protrusion}%
    {fontspec must be loaded before verse_protrusion.tex}%
    {Add \string\usepackage{fontspec} (and set a main font) before \string\input{verse_protrusion}.}%
}{}

% --- Tunable: partial left-protrusion factor (out of 1000 = one glyph width).
%     Edit ONLY this line to retune the hang; the \edef below propagates it
%     into all ten digit entries automatically.
\def\verseprot@factor{350}

% Build the entry list with the factor expanded in, NO spaces anywhere
% (catcode-10 spaces in the stored token list break \MT@pr@split@val --
%  see header note #4b).  Every newline is suppressed by a trailing `%'.
\edef\verseprot@list{%
"2070={\verseprot@factor,0},%
"00B9={\verseprot@factor,0},%
"00B2={\verseprot@factor,0},%
"00B3={\verseprot@factor,0},%
"2074={\verseprot@factor,0},%
"2075={\verseprot@factor,0},%
"2076={\verseprot@factor,0},%
"2077={\verseprot@factor,0},%
"2078={\verseprot@factor,0},%
"2079={\verseprot@factor,0}}

% --- Manual override (an edition may force the feature off):
\newif\ifverseprotrusion  \verseprotrusiontrue

% --- Global (default) microtype protrusion: left-only hang on the SUPERSCRIPT
%     glyphs that the `sups' feature substitutes for the digits 0-9 (see header
%     note #5). Keyed on the substituted Unicode codepoints, NOT on ASCII 0-9.
%     NO [name=...] and NO [context=...]: this MERGES into the protrusion list
%     microtype loads for the VerticalPosition=Superior font instance, which is
%     what actually binds the codepoints (header note #1). Safe globally because
%     these superscript glyphs only ever appear as verse numbers in this doc.
%     \verseprot@list is passed via \edef so the literal text (not the macro
%     name) is baked into the call before microtype reads it (header note #4).
% [load=default]: INHERIT microtype's built-in protrusion (quotes, hyphens,
% dashes, periods -- the hanging punctuation) and ADD the superscript-figure
% codepoints on top.  Without load=, a global \SetProtrusion REPLACES the
% font's whole protrusion vector, silently disabling all default hanging
% punctuation across the edition.
\edef\verseprot@do@protrusion{%
  \noexpand\SetProtrusion
    [ load = default ]%
    { font = */*/*/*/* }%
    {\verseprot@list}%
}%
\verseprot@do@protrusion

% Adapted for LaTeXGen: called in scripture's font hook AFTER selecting
% scripturefont. The sermon main font is intentionally a different family.
\newif\ifverseprot@use
\newcommand{\detectverseprotrusion}{%
  \setbox0=\hbox{{\addfontfeature{VerticalPosition=Superior}0}}%
  \setbox2=\hbox{0}%
  \ifdim\wd0=\wd2
    \verseprot@usefalse
  \else
    \ifverseprotrusion \verseprot@usetrue \else \verseprot@usefalse \fi
  \fi
}

% --- The formatter wired into scripture's verse/format.
\NewDocumentCommand\verseprotrudenum{m}{%
  \ifverseprot@use
    {\addfontfeature{VerticalPosition=Superior}#1}%
  \else
    \textsuperscript{#1}%
  \fi
}

\makeatother
.}%
}{}

% --- Tunable: partial left-protrusion factor (out of 1000 = one glyph width).
%     Edit ONLY this line to retune the hang; the \edef below propagates it
%     into all ten digit entries automatically.
\def\verseprot@factor{350}

% Build the entry list with the factor expanded in, NO spaces anywhere
% (catcode-10 spaces in the stored token list break \MT@pr@split@val --
%  see header note #4b).  Every newline is suppressed by a trailing `%'.
\edef\verseprot@list{%
"2070={\verseprot@factor,0},%
"00B9={\verseprot@factor,0},%
"00B2={\verseprot@factor,0},%
"00B3={\verseprot@factor,0},%
"2074={\verseprot@factor,0},%
"2075={\verseprot@factor,0},%
"2076={\verseprot@factor,0},%
"2077={\verseprot@factor,0},%
"2078={\verseprot@factor,0},%
"2079={\verseprot@factor,0}}

% --- Manual override (an edition may force the feature off):
\newif\ifverseprotrusion  \verseprotrusiontrue

% --- Global (default) microtype protrusion: left-only hang on the SUPERSCRIPT
%     glyphs that the `sups' feature substitutes for the digits 0-9 (see header
%     note #5). Keyed on the substituted Unicode codepoints, NOT on ASCII 0-9.
%     NO [name=...] and NO [context=...]: this MERGES into the protrusion list
%     microtype loads for the VerticalPosition=Superior font instance, which is
%     what actually binds the codepoints (header note #1). Safe globally because
%     these superscript glyphs only ever appear as verse numbers in this doc.
%     \verseprot@list is passed via \edef so the literal text (not the macro
%     name) is baked into the call before microtype reads it (header note #4).
% [load=default]: INHERIT microtype's built-in protrusion (quotes, hyphens,
% dashes, periods -- the hanging punctuation) and ADD the superscript-figure
% codepoints on top.  Without load=, a global \SetProtrusion REPLACES the
% font's whole protrusion vector, silently disabling all default hanging
% punctuation across the edition.
\edef\verseprot@do@protrusion{%
  \noexpand\SetProtrusion
    [ load = default ]%
    { font = */*/*/*/* }%
    {\verseprot@list}%
}%
\verseprot@do@protrusion

% Adapted for LaTeXGen: called in scripture's font hook AFTER selecting
% scripturefont. The sermon main font is intentionally a different family.
\newif\ifverseprot@use
\newcommand{\detectverseprotrusion}{%
  \setbox0=\hbox{{\addfontfeature{VerticalPosition=Superior}0}}%
  \setbox2=\hbox{0}%
  \ifdim\wd0=\wd2
    \verseprot@usefalse
  \else
    \ifverseprotrusion \verseprot@usetrue \else \verseprot@usefalse \fi
  \fi
}

% --- The formatter wired into scripture's verse/format.
\NewDocumentCommand\verseprotrudenum{m}{%
  \ifverseprot@use
    {\addfontfeature{VerticalPosition=Superior}#1}%
  \else
    \textsuperscript{#1}%
  \fi
}

\makeatother
.}%
}{}
\@ifundefined{addfontfeature}{%
  \PackageError{verse_protrusion}%
    {fontspec must be loaded before verse_protrusion.tex}%
    {Add \string\usepackage{fontspec} (and set a main font) before \string%% verse_protrusion.tex
%% Microtype protrusion for verse numbers (edition-portable).
%% Requires, loaded BEFORE this file: fontspec (main font set), microtype, scripture.
%% Effect: a verse number at the start of a line hangs a tuned partial amount
%% into the left margin. Mid-line numbers are unaffected (microtype only
%% protrudes line-edge glyphs). Font-agnostic: uses OpenType superior figures
%% of the CURRENT body font; falls back to \textsuperscript if unavailable.
%%
%% WORKING MECHANISM (the bits that are easy to get wrong):
%%   1. *** NO named context, NO \microtypecontext. *** Use a PLAIN, default
%%      \SetProtrusion (neither [name=...] nor [context=...]). microtype then
%%      MERGES these entries into whatever protrusion list it loads for any
%%      matching font, so the superior-figure instance gets the codepoints
%%      automatically. The earlier [context=verseprot] +
%%      \microtypecontext{protrusion=verseprot} approach was a NO-OP: the
%%      verbose log showed the context switch loaded an auto-generated EMPTY
%%      list named `<file>/<line>' and the named list never bound to the active
%%      superior font, so lpcode stayed 0 (measured: zero movement at 600 dpi).
%%      The global config is safe here because the superscript-codepoint glyphs
%%      (U+00B9 etc.) appear in this document ONLY as verse numbers (body prose
%%      never emits superior figures), so the entries are effectively scoped to
%%      verse numbers already. PROOF the binding now works: with the global
%%      config, `\the\lpcode\font 185' (U+00B9) reports a nonzero value on the
%%      VerticalPosition=Superior instance (was 0/unset before), and 600 dpi
%%      measurement shows the leading verse number's leftmost ink ~9 px left of
%%      where it sits with the feature off.
%%   2. Font set is the full 5-axis wildcard { font = */*/*/*/* } so it matches
%%      BOTH the body font instance AND the separate font instance that fontspec
%%      spins up for VerticalPosition=Superior (e.g. EBGaramond(1)).
%%   3. Each entry is char = {left,right} (ONE brace, comma-separated). Left-only
%%      hang -> {350,0}. Two separate braces {350}{0} is wrong syntax.
%%   4. To retune the factor, edit \verseprot@factor (one place). The \edef
%%      below expands it into all ten entries before passing to \SetProtrusion.
%%      TWO critical constraints on the expansion approach:
%%      (a) microtype stores \SetProtrusion's third arg via \gdef -- if we pass
%%          `{\verseprot@list}' as the arg, the macro TOKEN (not its expansion)
%%          is stored, and microtype's \MT@set@codes later sees \verseprot@list
%%          as a single opaque token, not as the hex-entry list. We therefore use
%%          `\edef\verseprot@do@protrusion{...{\verseprot@list}...}' to EXPAND the
%%          list inline and store the literal tokens in \verseprot@do@protrusion
%%          before \SetProtrusion is ever called.
%%      (b) microtype reads the \SetProtrusion list with \catcode` '=9 (spaces
%%          ignored), but our \edef runs with normal catcodes where space is
%%          catcode 10. If the expanded list contains catcode-10 space tokens,
%%          \MT@pr@split@val (which reads `left,right\relax') sees a leading space
%%          before `{350,0}' and can't find the `,\relax' terminator -> runaway.
%%          The fix: write the list with NO spaces (all newlines suppressed by `%'
%%          and no space around `=' or `{'). microtype's own parser ignores the
%%          lack of spaces; it just needs the correct token sequence.
%%   5. *** THE OTHER KEY FIX ***  microtype (luatex) keys protrusion by the codepoint
%%      of the glyph AS IT APPEARS IN THE NODE LIST, i.e. AFTER GSUB shaping.
%%      VerticalPosition=Superior runs the OpenType `sups' feature, which in
%%      EB Garamond substitutes each ASCII digit for the glyph at its Unicode
%%      SUPERSCRIPT codepoint:
%%          0->U+2070  1->U+00B9  2->U+00B2  3->U+00B3
%%          4..9 -> U+2074..U+2079
%%      Therefore the list MUST be keyed on those superscript codepoints, NOT on
%%      the ASCII digits 0-9. (Keying 0-9 = U+0030..U+0039 was the bug: those
%%      input glyphs never reach the line because sups replaces them, so the
%%      lpcode never bound to the rendered superior glyph -> ZERO protrusion,
%%      verified by 600 dpi pixel measurement and by dumping the substituted
%%      glyph names with fontTools: three -> uni00B3, etc.)
%%   6. microtype protrudes GLYPH nodes, not boxes. The substituted superior
%%      figures are real glyphs and protrude; \textsuperscript (a raised box)
%%      does not -- hence the OpenType path is required for the hang, with
%%      \textsuperscript only as the non-protruding fallback when the font lacks
%%      superior figures.
%% Verified by measurement (tests/verse_protrusion_visual_check.tex, 600 dpi):
%% with this setup the leading verse number's leftmost ink sits measurably LEFT
%% of the column rule, while mid-line numbers stay inline; feature ON vs OFF
%% now differs at the left edge.
%%
%% PORTABILITY NOTE: the superscript-codepoint mapping above is what EB Garamond
%% (and most fonts following the Unicode/AGL convention) emit for `sups'. A font
%% whose superior figures live at private-use or `.sups'-named glyph slots with
%% no Unicode value would need different keys; in that case microtype cannot
%% address them by codepoint and the \textsuperscript fallback (no protrusion)
%% is the safe outcome.
\makeatletter

% --- Precondition guards: raise a clear error if required packages are absent.
%     Without microtype (\SetProtrusion) or fontspec (\addfontfeature) the file
%     would fail later with a cryptic "Undefined control sequence" inside an \edef.
\@ifundefined{SetProtrusion}{%
  \PackageError{verse_protrusion}%
    {microtype must be loaded before verse_protrusion.tex}%
    {Add \string\usepackage{microtype} (and fontspec) before \string%% verse_protrusion.tex
%% Microtype protrusion for verse numbers (edition-portable).
%% Requires, loaded BEFORE this file: fontspec (main font set), microtype, scripture.
%% Effect: a verse number at the start of a line hangs a tuned partial amount
%% into the left margin. Mid-line numbers are unaffected (microtype only
%% protrudes line-edge glyphs). Font-agnostic: uses OpenType superior figures
%% of the CURRENT body font; falls back to \textsuperscript if unavailable.
%%
%% WORKING MECHANISM (the bits that are easy to get wrong):
%%   1. *** NO named context, NO \microtypecontext. *** Use a PLAIN, default
%%      \SetProtrusion (neither [name=...] nor [context=...]). microtype then
%%      MERGES these entries into whatever protrusion list it loads for any
%%      matching font, so the superior-figure instance gets the codepoints
%%      automatically. The earlier [context=verseprot] +
%%      \microtypecontext{protrusion=verseprot} approach was a NO-OP: the
%%      verbose log showed the context switch loaded an auto-generated EMPTY
%%      list named `<file>/<line>' and the named list never bound to the active
%%      superior font, so lpcode stayed 0 (measured: zero movement at 600 dpi).
%%      The global config is safe here because the superscript-codepoint glyphs
%%      (U+00B9 etc.) appear in this document ONLY as verse numbers (body prose
%%      never emits superior figures), so the entries are effectively scoped to
%%      verse numbers already. PROOF the binding now works: with the global
%%      config, `\the\lpcode\font 185' (U+00B9) reports a nonzero value on the
%%      VerticalPosition=Superior instance (was 0/unset before), and 600 dpi
%%      measurement shows the leading verse number's leftmost ink ~9 px left of
%%      where it sits with the feature off.
%%   2. Font set is the full 5-axis wildcard { font = */*/*/*/* } so it matches
%%      BOTH the body font instance AND the separate font instance that fontspec
%%      spins up for VerticalPosition=Superior (e.g. EBGaramond(1)).
%%   3. Each entry is char = {left,right} (ONE brace, comma-separated). Left-only
%%      hang -> {350,0}. Two separate braces {350}{0} is wrong syntax.
%%   4. To retune the factor, edit \verseprot@factor (one place). The \edef
%%      below expands it into all ten entries before passing to \SetProtrusion.
%%      TWO critical constraints on the expansion approach:
%%      (a) microtype stores \SetProtrusion's third arg via \gdef -- if we pass
%%          `{\verseprot@list}' as the arg, the macro TOKEN (not its expansion)
%%          is stored, and microtype's \MT@set@codes later sees \verseprot@list
%%          as a single opaque token, not as the hex-entry list. We therefore use
%%          `\edef\verseprot@do@protrusion{...{\verseprot@list}...}' to EXPAND the
%%          list inline and store the literal tokens in \verseprot@do@protrusion
%%          before \SetProtrusion is ever called.
%%      (b) microtype reads the \SetProtrusion list with \catcode` '=9 (spaces
%%          ignored), but our \edef runs with normal catcodes where space is
%%          catcode 10. If the expanded list contains catcode-10 space tokens,
%%          \MT@pr@split@val (which reads `left,right\relax') sees a leading space
%%          before `{350,0}' and can't find the `,\relax' terminator -> runaway.
%%          The fix: write the list with NO spaces (all newlines suppressed by `%'
%%          and no space around `=' or `{'). microtype's own parser ignores the
%%          lack of spaces; it just needs the correct token sequence.
%%   5. *** THE OTHER KEY FIX ***  microtype (luatex) keys protrusion by the codepoint
%%      of the glyph AS IT APPEARS IN THE NODE LIST, i.e. AFTER GSUB shaping.
%%      VerticalPosition=Superior runs the OpenType `sups' feature, which in
%%      EB Garamond substitutes each ASCII digit for the glyph at its Unicode
%%      SUPERSCRIPT codepoint:
%%          0->U+2070  1->U+00B9  2->U+00B2  3->U+00B3
%%          4..9 -> U+2074..U+2079
%%      Therefore the list MUST be keyed on those superscript codepoints, NOT on
%%      the ASCII digits 0-9. (Keying 0-9 = U+0030..U+0039 was the bug: those
%%      input glyphs never reach the line because sups replaces them, so the
%%      lpcode never bound to the rendered superior glyph -> ZERO protrusion,
%%      verified by 600 dpi pixel measurement and by dumping the substituted
%%      glyph names with fontTools: three -> uni00B3, etc.)
%%   6. microtype protrudes GLYPH nodes, not boxes. The substituted superior
%%      figures are real glyphs and protrude; \textsuperscript (a raised box)
%%      does not -- hence the OpenType path is required for the hang, with
%%      \textsuperscript only as the non-protruding fallback when the font lacks
%%      superior figures.
%% Verified by measurement (tests/verse_protrusion_visual_check.tex, 600 dpi):
%% with this setup the leading verse number's leftmost ink sits measurably LEFT
%% of the column rule, while mid-line numbers stay inline; feature ON vs OFF
%% now differs at the left edge.
%%
%% PORTABILITY NOTE: the superscript-codepoint mapping above is what EB Garamond
%% (and most fonts following the Unicode/AGL convention) emit for `sups'. A font
%% whose superior figures live at private-use or `.sups'-named glyph slots with
%% no Unicode value would need different keys; in that case microtype cannot
%% address them by codepoint and the \textsuperscript fallback (no protrusion)
%% is the safe outcome.
\makeatletter

% --- Precondition guards: raise a clear error if required packages are absent.
%     Without microtype (\SetProtrusion) or fontspec (\addfontfeature) the file
%     would fail later with a cryptic "Undefined control sequence" inside an \edef.
\@ifundefined{SetProtrusion}{%
  \PackageError{verse_protrusion}%
    {microtype must be loaded before verse_protrusion.tex}%
    {Add \string\usepackage{microtype} (and fontspec) before \string\input{verse_protrusion}.}%
}{}
\@ifundefined{addfontfeature}{%
  \PackageError{verse_protrusion}%
    {fontspec must be loaded before verse_protrusion.tex}%
    {Add \string\usepackage{fontspec} (and set a main font) before \string\input{verse_protrusion}.}%
}{}

% --- Tunable: partial left-protrusion factor (out of 1000 = one glyph width).
%     Edit ONLY this line to retune the hang; the \edef below propagates it
%     into all ten digit entries automatically.
\def\verseprot@factor{350}

% Build the entry list with the factor expanded in, NO spaces anywhere
% (catcode-10 spaces in the stored token list break \MT@pr@split@val --
%  see header note #4b).  Every newline is suppressed by a trailing `%'.
\edef\verseprot@list{%
"2070={\verseprot@factor,0},%
"00B9={\verseprot@factor,0},%
"00B2={\verseprot@factor,0},%
"00B3={\verseprot@factor,0},%
"2074={\verseprot@factor,0},%
"2075={\verseprot@factor,0},%
"2076={\verseprot@factor,0},%
"2077={\verseprot@factor,0},%
"2078={\verseprot@factor,0},%
"2079={\verseprot@factor,0}}

% --- Manual override (an edition may force the feature off):
\newif\ifverseprotrusion  \verseprotrusiontrue

% --- Global (default) microtype protrusion: left-only hang on the SUPERSCRIPT
%     glyphs that the `sups' feature substitutes for the digits 0-9 (see header
%     note #5). Keyed on the substituted Unicode codepoints, NOT on ASCII 0-9.
%     NO [name=...] and NO [context=...]: this MERGES into the protrusion list
%     microtype loads for the VerticalPosition=Superior font instance, which is
%     what actually binds the codepoints (header note #1). Safe globally because
%     these superscript glyphs only ever appear as verse numbers in this doc.
%     \verseprot@list is passed via \edef so the literal text (not the macro
%     name) is baked into the call before microtype reads it (header note #4).
% [load=default]: INHERIT microtype's built-in protrusion (quotes, hyphens,
% dashes, periods -- the hanging punctuation) and ADD the superscript-figure
% codepoints on top.  Without load=, a global \SetProtrusion REPLACES the
% font's whole protrusion vector, silently disabling all default hanging
% punctuation across the edition.
\edef\verseprot@do@protrusion{%
  \noexpand\SetProtrusion
    [ load = default ]%
    { font = */*/*/*/* }%
    {\verseprot@list}%
}%
\verseprot@do@protrusion

% Adapted for LaTeXGen: called in scripture's font hook AFTER selecting
% scripturefont. The sermon main font is intentionally a different family.
\newif\ifverseprot@use
\newcommand{\detectverseprotrusion}{%
  \setbox0=\hbox{{\addfontfeature{VerticalPosition=Superior}0}}%
  \setbox2=\hbox{0}%
  \ifdim\wd0=\wd2
    \verseprot@usefalse
  \else
    \ifverseprotrusion \verseprot@usetrue \else \verseprot@usefalse \fi
  \fi
}

% --- The formatter wired into scripture's verse/format.
\NewDocumentCommand\verseprotrudenum{m}{%
  \ifverseprot@use
    {\addfontfeature{VerticalPosition=Superior}#1}%
  \else
    \textsuperscript{#1}%
  \fi
}

\makeatother
.}%
}{}
\@ifundefined{addfontfeature}{%
  \PackageError{verse_protrusion}%
    {fontspec must be loaded before verse_protrusion.tex}%
    {Add \string\usepackage{fontspec} (and set a main font) before \string%% verse_protrusion.tex
%% Microtype protrusion for verse numbers (edition-portable).
%% Requires, loaded BEFORE this file: fontspec (main font set), microtype, scripture.
%% Effect: a verse number at the start of a line hangs a tuned partial amount
%% into the left margin. Mid-line numbers are unaffected (microtype only
%% protrudes line-edge glyphs). Font-agnostic: uses OpenType superior figures
%% of the CURRENT body font; falls back to \textsuperscript if unavailable.
%%
%% WORKING MECHANISM (the bits that are easy to get wrong):
%%   1. *** NO named context, NO \microtypecontext. *** Use a PLAIN, default
%%      \SetProtrusion (neither [name=...] nor [context=...]). microtype then
%%      MERGES these entries into whatever protrusion list it loads for any
%%      matching font, so the superior-figure instance gets the codepoints
%%      automatically. The earlier [context=verseprot] +
%%      \microtypecontext{protrusion=verseprot} approach was a NO-OP: the
%%      verbose log showed the context switch loaded an auto-generated EMPTY
%%      list named `<file>/<line>' and the named list never bound to the active
%%      superior font, so lpcode stayed 0 (measured: zero movement at 600 dpi).
%%      The global config is safe here because the superscript-codepoint glyphs
%%      (U+00B9 etc.) appear in this document ONLY as verse numbers (body prose
%%      never emits superior figures), so the entries are effectively scoped to
%%      verse numbers already. PROOF the binding now works: with the global
%%      config, `\the\lpcode\font 185' (U+00B9) reports a nonzero value on the
%%      VerticalPosition=Superior instance (was 0/unset before), and 600 dpi
%%      measurement shows the leading verse number's leftmost ink ~9 px left of
%%      where it sits with the feature off.
%%   2. Font set is the full 5-axis wildcard { font = */*/*/*/* } so it matches
%%      BOTH the body font instance AND the separate font instance that fontspec
%%      spins up for VerticalPosition=Superior (e.g. EBGaramond(1)).
%%   3. Each entry is char = {left,right} (ONE brace, comma-separated). Left-only
%%      hang -> {350,0}. Two separate braces {350}{0} is wrong syntax.
%%   4. To retune the factor, edit \verseprot@factor (one place). The \edef
%%      below expands it into all ten entries before passing to \SetProtrusion.
%%      TWO critical constraints on the expansion approach:
%%      (a) microtype stores \SetProtrusion's third arg via \gdef -- if we pass
%%          `{\verseprot@list}' as the arg, the macro TOKEN (not its expansion)
%%          is stored, and microtype's \MT@set@codes later sees \verseprot@list
%%          as a single opaque token, not as the hex-entry list. We therefore use
%%          `\edef\verseprot@do@protrusion{...{\verseprot@list}...}' to EXPAND the
%%          list inline and store the literal tokens in \verseprot@do@protrusion
%%          before \SetProtrusion is ever called.
%%      (b) microtype reads the \SetProtrusion list with \catcode` '=9 (spaces
%%          ignored), but our \edef runs with normal catcodes where space is
%%          catcode 10. If the expanded list contains catcode-10 space tokens,
%%          \MT@pr@split@val (which reads `left,right\relax') sees a leading space
%%          before `{350,0}' and can't find the `,\relax' terminator -> runaway.
%%          The fix: write the list with NO spaces (all newlines suppressed by `%'
%%          and no space around `=' or `{'). microtype's own parser ignores the
%%          lack of spaces; it just needs the correct token sequence.
%%   5. *** THE OTHER KEY FIX ***  microtype (luatex) keys protrusion by the codepoint
%%      of the glyph AS IT APPEARS IN THE NODE LIST, i.e. AFTER GSUB shaping.
%%      VerticalPosition=Superior runs the OpenType `sups' feature, which in
%%      EB Garamond substitutes each ASCII digit for the glyph at its Unicode
%%      SUPERSCRIPT codepoint:
%%          0->U+2070  1->U+00B9  2->U+00B2  3->U+00B3
%%          4..9 -> U+2074..U+2079
%%      Therefore the list MUST be keyed on those superscript codepoints, NOT on
%%      the ASCII digits 0-9. (Keying 0-9 = U+0030..U+0039 was the bug: those
%%      input glyphs never reach the line because sups replaces them, so the
%%      lpcode never bound to the rendered superior glyph -> ZERO protrusion,
%%      verified by 600 dpi pixel measurement and by dumping the substituted
%%      glyph names with fontTools: three -> uni00B3, etc.)
%%   6. microtype protrudes GLYPH nodes, not boxes. The substituted superior
%%      figures are real glyphs and protrude; \textsuperscript (a raised box)
%%      does not -- hence the OpenType path is required for the hang, with
%%      \textsuperscript only as the non-protruding fallback when the font lacks
%%      superior figures.
%% Verified by measurement (tests/verse_protrusion_visual_check.tex, 600 dpi):
%% with this setup the leading verse number's leftmost ink sits measurably LEFT
%% of the column rule, while mid-line numbers stay inline; feature ON vs OFF
%% now differs at the left edge.
%%
%% PORTABILITY NOTE: the superscript-codepoint mapping above is what EB Garamond
%% (and most fonts following the Unicode/AGL convention) emit for `sups'. A font
%% whose superior figures live at private-use or `.sups'-named glyph slots with
%% no Unicode value would need different keys; in that case microtype cannot
%% address them by codepoint and the \textsuperscript fallback (no protrusion)
%% is the safe outcome.
\makeatletter

% --- Precondition guards: raise a clear error if required packages are absent.
%     Without microtype (\SetProtrusion) or fontspec (\addfontfeature) the file
%     would fail later with a cryptic "Undefined control sequence" inside an \edef.
\@ifundefined{SetProtrusion}{%
  \PackageError{verse_protrusion}%
    {microtype must be loaded before verse_protrusion.tex}%
    {Add \string\usepackage{microtype} (and fontspec) before \string\input{verse_protrusion}.}%
}{}
\@ifundefined{addfontfeature}{%
  \PackageError{verse_protrusion}%
    {fontspec must be loaded before verse_protrusion.tex}%
    {Add \string\usepackage{fontspec} (and set a main font) before \string\input{verse_protrusion}.}%
}{}

% --- Tunable: partial left-protrusion factor (out of 1000 = one glyph width).
%     Edit ONLY this line to retune the hang; the \edef below propagates it
%     into all ten digit entries automatically.
\def\verseprot@factor{350}

% Build the entry list with the factor expanded in, NO spaces anywhere
% (catcode-10 spaces in the stored token list break \MT@pr@split@val --
%  see header note #4b).  Every newline is suppressed by a trailing `%'.
\edef\verseprot@list{%
"2070={\verseprot@factor,0},%
"00B9={\verseprot@factor,0},%
"00B2={\verseprot@factor,0},%
"00B3={\verseprot@factor,0},%
"2074={\verseprot@factor,0},%
"2075={\verseprot@factor,0},%
"2076={\verseprot@factor,0},%
"2077={\verseprot@factor,0},%
"2078={\verseprot@factor,0},%
"2079={\verseprot@factor,0}}

% --- Manual override (an edition may force the feature off):
\newif\ifverseprotrusion  \verseprotrusiontrue

% --- Global (default) microtype protrusion: left-only hang on the SUPERSCRIPT
%     glyphs that the `sups' feature substitutes for the digits 0-9 (see header
%     note #5). Keyed on the substituted Unicode codepoints, NOT on ASCII 0-9.
%     NO [name=...] and NO [context=...]: this MERGES into the protrusion list
%     microtype loads for the VerticalPosition=Superior font instance, which is
%     what actually binds the codepoints (header note #1). Safe globally because
%     these superscript glyphs only ever appear as verse numbers in this doc.
%     \verseprot@list is passed via \edef so the literal text (not the macro
%     name) is baked into the call before microtype reads it (header note #4).
% [load=default]: INHERIT microtype's built-in protrusion (quotes, hyphens,
% dashes, periods -- the hanging punctuation) and ADD the superscript-figure
% codepoints on top.  Without load=, a global \SetProtrusion REPLACES the
% font's whole protrusion vector, silently disabling all default hanging
% punctuation across the edition.
\edef\verseprot@do@protrusion{%
  \noexpand\SetProtrusion
    [ load = default ]%
    { font = */*/*/*/* }%
    {\verseprot@list}%
}%
\verseprot@do@protrusion

% Adapted for LaTeXGen: called in scripture's font hook AFTER selecting
% scripturefont. The sermon main font is intentionally a different family.
\newif\ifverseprot@use
\newcommand{\detectverseprotrusion}{%
  \setbox0=\hbox{{\addfontfeature{VerticalPosition=Superior}0}}%
  \setbox2=\hbox{0}%
  \ifdim\wd0=\wd2
    \verseprot@usefalse
  \else
    \ifverseprotrusion \verseprot@usetrue \else \verseprot@usefalse \fi
  \fi
}

% --- The formatter wired into scripture's verse/format.
\NewDocumentCommand\verseprotrudenum{m}{%
  \ifverseprot@use
    {\addfontfeature{VerticalPosition=Superior}#1}%
  \else
    \textsuperscript{#1}%
  \fi
}

\makeatother
.}%
}{}

% --- Tunable: partial left-protrusion factor (out of 1000 = one glyph width).
%     Edit ONLY this line to retune the hang; the \edef below propagates it
%     into all ten digit entries automatically.
\def\verseprot@factor{350}

% Build the entry list with the factor expanded in, NO spaces anywhere
% (catcode-10 spaces in the stored token list break \MT@pr@split@val --
%  see header note #4b).  Every newline is suppressed by a trailing `%'.
\edef\verseprot@list{%
"2070={\verseprot@factor,0},%
"00B9={\verseprot@factor,0},%
"00B2={\verseprot@factor,0},%
"00B3={\verseprot@factor,0},%
"2074={\verseprot@factor,0},%
"2075={\verseprot@factor,0},%
"2076={\verseprot@factor,0},%
"2077={\verseprot@factor,0},%
"2078={\verseprot@factor,0},%
"2079={\verseprot@factor,0}}

% --- Manual override (an edition may force the feature off):
\newif\ifverseprotrusion  \verseprotrusiontrue

% --- Global (default) microtype protrusion: left-only hang on the SUPERSCRIPT
%     glyphs that the `sups' feature substitutes for the digits 0-9 (see header
%     note #5). Keyed on the substituted Unicode codepoints, NOT on ASCII 0-9.
%     NO [name=...] and NO [context=...]: this MERGES into the protrusion list
%     microtype loads for the VerticalPosition=Superior font instance, which is
%     what actually binds the codepoints (header note #1). Safe globally because
%     these superscript glyphs only ever appear as verse numbers in this doc.
%     \verseprot@list is passed via \edef so the literal text (not the macro
%     name) is baked into the call before microtype reads it (header note #4).
% [load=default]: INHERIT microtype's built-in protrusion (quotes, hyphens,
% dashes, periods -- the hanging punctuation) and ADD the superscript-figure
% codepoints on top.  Without load=, a global \SetProtrusion REPLACES the
% font's whole protrusion vector, silently disabling all default hanging
% punctuation across the edition.
\edef\verseprot@do@protrusion{%
  \noexpand\SetProtrusion
    [ load = default ]%
    { font = */*/*/*/* }%
    {\verseprot@list}%
}%
\verseprot@do@protrusion

% Adapted for LaTeXGen: called in scripture's font hook AFTER selecting
% scripturefont. The sermon main font is intentionally a different family.
\newif\ifverseprot@use
\newcommand{\detectverseprotrusion}{%
  \setbox0=\hbox{{\addfontfeature{VerticalPosition=Superior}0}}%
  \setbox2=\hbox{0}%
  \ifdim\wd0=\wd2
    \verseprot@usefalse
  \else
    \ifverseprotrusion \verseprot@usetrue \else \verseprot@usefalse \fi
  \fi
}

% --- The formatter wired into scripture's verse/format.
\NewDocumentCommand\verseprotrudenum{m}{%
  \ifverseprot@use
    {\addfontfeature{VerticalPosition=Superior}#1}%
  \else
    \textsuperscript{#1}%
  \fi
}

\makeatother
.}%
}{}

% --- Tunable: partial left-protrusion factor (out of 1000 = one glyph width).
%     Edit ONLY this line to retune the hang; the \edef below propagates it
%     into all ten digit entries automatically.
\def\verseprot@factor{350}

% Build the entry list with the factor expanded in, NO spaces anywhere
% (catcode-10 spaces in the stored token list break \MT@pr@split@val --
%  see header note #4b).  Every newline is suppressed by a trailing `%'.
\edef\verseprot@list{%
"2070={\verseprot@factor,0},%
"00B9={\verseprot@factor,0},%
"00B2={\verseprot@factor,0},%
"00B3={\verseprot@factor,0},%
"2074={\verseprot@factor,0},%
"2075={\verseprot@factor,0},%
"2076={\verseprot@factor,0},%
"2077={\verseprot@factor,0},%
"2078={\verseprot@factor,0},%
"2079={\verseprot@factor,0}}

% --- Manual override (an edition may force the feature off):
\newif\ifverseprotrusion  \verseprotrusiontrue

% --- Global (default) microtype protrusion: left-only hang on the SUPERSCRIPT
%     glyphs that the `sups' feature substitutes for the digits 0-9 (see header
%     note #5). Keyed on the substituted Unicode codepoints, NOT on ASCII 0-9.
%     NO [name=...] and NO [context=...]: this MERGES into the protrusion list
%     microtype loads for the VerticalPosition=Superior font instance, which is
%     what actually binds the codepoints (header note #1). Safe globally because
%     these superscript glyphs only ever appear as verse numbers in this doc.
%     \verseprot@list is passed via \edef so the literal text (not the macro
%     name) is baked into the call before microtype reads it (header note #4).
% [load=default]: INHERIT microtype's built-in protrusion (quotes, hyphens,
% dashes, periods -- the hanging punctuation) and ADD the superscript-figure
% codepoints on top.  Without load=, a global \SetProtrusion REPLACES the
% font's whole protrusion vector, silently disabling all default hanging
% punctuation across the edition.
\edef\verseprot@do@protrusion{%
  \noexpand\SetProtrusion
    [ load = default ]%
    { font = */*/*/*/* }%
    {\verseprot@list}%
}%
\verseprot@do@protrusion

% Adapted for LaTeXGen: called in scripture's font hook AFTER selecting
% scripturefont. The sermon main font is intentionally a different family.
\newif\ifverseprot@use
\newcommand{\detectverseprotrusion}{%
  \setbox0=\hbox{{\addfontfeature{VerticalPosition=Superior}0}}%
  \setbox2=\hbox{0}%
  \ifdim\wd0=\wd2
    \verseprot@usefalse
  \else
    \ifverseprotrusion \verseprot@usetrue \else \verseprot@usefalse \fi
  \fi
}

% --- The formatter wired into scripture's verse/format.
\NewDocumentCommand\verseprotrudenum{m}{%
  \ifverseprot@use
    {\addfontfeature{VerticalPosition=Superior}#1}%
  \else
    \textsuperscript{#1}%
  \fi
}

\makeatother
.}%
}{}

% --- Tunable: partial left-protrusion factor (out of 1000 = one glyph width).
%     Edit ONLY this line to retune the hang; the \edef below propagates it
%     into all ten digit entries automatically.
\def\verseprot@factor{350}

% Build the entry list with the factor expanded in, NO spaces anywhere
% (catcode-10 spaces in the stored token list break \MT@pr@split@val --
%  see header note #4b).  Every newline is suppressed by a trailing `%'.
\edef\verseprot@list{%
"2070={\verseprot@factor,0},%
"00B9={\verseprot@factor,0},%
"00B2={\verseprot@factor,0},%
"00B3={\verseprot@factor,0},%
"2074={\verseprot@factor,0},%
"2075={\verseprot@factor,0},%
"2076={\verseprot@factor,0},%
"2077={\verseprot@factor,0},%
"2078={\verseprot@factor,0},%
"2079={\verseprot@factor,0}}

% --- Manual override (an edition may force the feature off):
\newif\ifverseprotrusion  \verseprotrusiontrue

% --- Global (default) microtype protrusion: left-only hang on the SUPERSCRIPT
%     glyphs that the `sups' feature substitutes for the digits 0-9 (see header
%     note #5). Keyed on the substituted Unicode codepoints, NOT on ASCII 0-9.
%     NO [name=...] and NO [context=...]: this MERGES into the protrusion list
%     microtype loads for the VerticalPosition=Superior font instance, which is
%     what actually binds the codepoints (header note #1). Safe globally because
%     these superscript glyphs only ever appear as verse numbers in this doc.
%     \verseprot@list is passed via \edef so the literal text (not the macro
%     name) is baked into the call before microtype reads it (header note #4).
% [load=default]: INHERIT microtype's built-in protrusion (quotes, hyphens,
% dashes, periods -- the hanging punctuation) and ADD the superscript-figure
% codepoints on top.  Without load=, a global \SetProtrusion REPLACES the
% font's whole protrusion vector, silently disabling all default hanging
% punctuation across the edition.
\edef\verseprot@do@protrusion{%
  \noexpand\SetProtrusion
    [ load = default ]%
    { font = */*/*/*/* }%
    {\verseprot@list}%
}%
\verseprot@do@protrusion

% Adapted for LaTeXGen: called in scripture's font hook AFTER selecting
% scripturefont. The sermon main font is intentionally a different family.
\newif\ifverseprot@use
\newcommand{\detectverseprotrusion}{%
  \setbox0=\hbox{{\addfontfeature{VerticalPosition=Superior}0}}%
  \setbox2=\hbox{0}%
  \ifdim\wd0=\wd2
    \verseprot@usefalse
  \else
    \ifverseprotrusion \verseprot@usetrue \else \verseprot@usefalse \fi
  \fi
}

% --- The formatter wired into scripture's verse/format.
\NewDocumentCommand\verseprotrudenum{m}{%
  \ifverseprot@use
    {\addfontfeature{VerticalPosition=Superior}#1}%
  \else
    \textsuperscript{#1}%
  \fi
}

\makeatother
